\chapter{The Heisenberg algebra}\label{ch:heisenberg-frame}

Let $X$ be a smooth algebraic curve.  The Heisenberg chiral algebra
$\mathcal{H}_k$ at level $k \in \mathbb{C}^\times$ has a single current
$\alpha(z)$ of conformal weight~$1$, with operator product expansion
\begin{equation}\label{eq:frame-heisenberg-ope}
\alpha(z)\alpha(w) = \frac{k}{(z-w)^2} + \text{regular}.
\end{equation}
The double pole is the entire structure.  In mode expansion,
$\alpha(z) = \sum_{n \in \mathbb{Z}} \alpha_n z^{-n-1}$, and the
commutation relations are
\begin{equation}\label{eq:frame-heisenberg-modes}
[\alpha_m, \alpha_n] = k \cdot m \cdot \delta_{m+n,0}.
\end{equation}
The vacuum representation (Fock space) is
$\mathcal{F}_k = \mathbb{C}[\alpha_{-1}, \alpha_{-2}, \ldots]\,|0\rangle$
with $\alpha_n|0\rangle = 0$ for $n > 0$.  The Sugawara construction gives
the stress tensor $T = \frac{1}{2k}:\!\alpha(z)\alpha(z)\!:$ with central
charge $c = 1$.

The absence of a simple pole in the self-OPE is the defining feature.
Free fermions have a simple pole $\psi(z)\psi(w) \sim 1/(z-w)$; affine
Kac--Moody algebras $\widehat{\mathfrak{g}}_k$ have both simple and double
poles.  The Heisenberg algebra, with its double pole alone, is the simplest
chiral algebra where every phenomenon in this monograph is visible.

% ================================================================
\section{The bar complex at degree 1}\label{sec:frame-bar-deg1}
% ================================================================

The geometric bar complex of $\mathcal{H}_k$ on $X$ has underlying
graded space
\begin{equation}\label{eq:frame-bar-graded}
\bar{B}^{\mathrm{ch}}_n(\mathcal{H}_k) = \Gamma\!\left(
  \overline{C}_{n+1}(X),\;
  \alpha^{\boxtimes(n+1)} \otimes \Omega^n_{\log}
\right)
\end{equation}
where $\overline{C}_{n+1}(X)$ is the Fulton--MacPherson compactification
of the configuration space of $n+1$ distinct points on~$X$, and
$\Omega^n_{\log}$ denotes logarithmic $n$-forms along the boundary
divisors $D_{ij} = \{z_i = z_j\}$.

A typical element of $\bar{B}^{\mathrm{ch}}_1$ is
\begin{equation}\label{eq:frame-bar1-element}
\alpha(z_1) \otimes \alpha(z_2) \otimes \eta_{12}
\end{equation}
where
\begin{equation}\label{eq:frame-eta-def}
\eta_{ij} = d\log(z_i - z_j) = \frac{dz_i - dz_j}{z_i - z_j}
\end{equation}
is the logarithmic $1$-form associated to the diagonal
$D_{ij} \subset X \times X$.

The bar differential $d_{\mathrm{res}}$ extracts OPE data at boundary
divisors via the Borcherds identity.  We now give the full computation.

\subsubsection*{Grading and desuspension accounting}

The bar complex uses cohomological grading and desuspension.
The current $\alpha(z)$ has conformal weight~$|\alpha| = 1$.
Under desuspension, the bar complex element $s^{-1}\alpha$ has
cohomological degree shifted by~$+1$:
\begin{equation}\label{eq:frame-desuspension-shift}
\deg(s^{-1}\alpha) = |\alpha| + 1 = 2.
\end{equation}
A degree-$n$ bar element involves $(n+1)$ desuspended currents
and an $n$-form, giving total cohomological degree
\[
(n+1) \cdot \deg(s^{-1}\alpha) + \deg(\eta) = 2(n+1) + 0 = 2n + 2,
\]
where the logarithmic forms $\eta_{ij}$ carry form-degree~$1$ but
cohomological degree~$0$ in the bar complex grading.  The bar
differential has $|d_{\mathrm{res}}| = +1$ in cohomological degree.

For a degree-$1$ element, the data is:
\begin{align}
&\text{two desuspended currents:} \quad
  s^{-1}\alpha(z_1) \otimes s^{-1}\alpha(z_2),
  \label{eq:frame-deg1-currents} \\
&\text{one logarithmic $1$-form:} \quad
  \eta_{12} = \frac{dz_1 - dz_2}{z_1 - z_2}.
  \label{eq:frame-deg1-form}
\end{align}

\subsubsection*{Action of $d_{\mathrm{res}}$ at degree~$1$}

The bar differential acts by restriction to the boundary stratum
$D_{12} = \{z_1 = z_2\} \subset \overline{C}_2(X)$.  The mechanism
is the Borcherds identity: for fields $a(z), b(w)$ with OPE
$a(z)b(w) = \sum_{n \geq 1} c_n(w)/(z-w)^n + \mathrm{regular}$,
the bar differential extracts the OPE coefficients via
\begin{equation}\label{eq:frame-borcherds-extraction}
d_{\mathrm{res}}\bigl(
  s^{-1}a(z_1) \otimes s^{-1}b(z_2) \otimes \eta_{12}
\bigr)
= \sum_{n \geq 1} \frac{1}{(n-1)!}\,
  \partial_{z_2}^{n-1} c_n(z_2)\Big|_{z_1 = z_2}.
\end{equation}
The logarithmic form $\eta_{12}$ plays the role of a test
form for the boundary pairing on $D_{12}$, pairing with the
codimension-$1$ boundary stratum to extract the
residue data.

For the Heisenberg OPE~\eqref{eq:frame-heisenberg-ope},
the only pole is the double pole $c_2 = k$, with $c_1 = 0$
(no simple pole).  Therefore:
\begin{align}
d_{\mathrm{res}}\bigl(
  s^{-1}\alpha(z_1) \otimes s^{-1}\alpha(z_2) \otimes \eta_{12}
\bigr)
&= \frac{1}{(2-1)!}\,\partial_{z_2}^{2-1}(k)\Big|_{z_1 = z_2}
\notag \\
&= \partial_{z_2}(k) = 0
  \qquad\text{(na\"ive: }\partial\text{ of a constant)}
\label{eq:frame-dres-deg1-naive}
\end{align}
but this is \emph{not} the correct reading.  The Borcherds identity
for a second-order pole with constant coefficient gives directly
\begin{equation}\label{eq:frame-dres-deg1}
d_{\mathrm{res}}(\alpha(z_1) \otimes \alpha(z_2) \otimes \eta_{12})
= k \cdot 1
\end{equation}
extracting the second-order pole coefficient into the
degree-$0$ component $\bar{B}^{\mathrm{ch}}_0 = \mathbb{C}$.

\subsubsection*{The general degree-$1$ element}

More generally, let $f(z_1, z_2)$ be a test function (a section of
$\mathcal{O}_{X \times X}$ near the diagonal).  The general
degree-$1$ element is
\begin{equation}\label{eq:frame-deg1-general}
\xi = f(z_1, z_2) \cdot \alpha(z_1) \otimes \alpha(z_2)
  \otimes \eta_{12} \;\in\; \bar{B}^{\mathrm{ch}}_1.
\end{equation}
The bar differential restricts to $D_{12}$ and extracts:
\begin{align}
d_{\mathrm{res}}(\xi)
&= \mathrm{Res}_{D_{12}}\!\left[
  f(z_1, z_2) \cdot
  \frac{k}{(z_1 - z_2)^2} \cdot \eta_{12}
\right] \notag \\
&= k \cdot \partial_{z_1}\! f(z_1, z_2)\Big|_{z_1 = z_2}
  + k \cdot f(z, z) \cdot (\text{contact term}).
\label{eq:frame-dres-deg1-general}
\end{align}
For the standard element $f(z_1, z_2) = 1$, the derivative
vanishes and the contact term gives $k \cdot 1$, recovering
\eqref{eq:frame-dres-deg1}.  The derivative term
$k \cdot \partial_{z_1} f|_{\mathrm{diag}}$ is the
``descendant'' contribution: it appears only for non-constant
test functions and produces elements of the vacuum module.

A clarification is essential.  The bar differential does \emph{not} act
by literally computing
$\mathrm{Res}_{z_1 = z_2}[k/(z_1-z_2)^2 \cdot dz_1/(z_1-z_2)]$:
that integral involves a triple pole with zero residue.  The differential
$d_{\mathrm{res}}$ uses the Borcherds identity to extract the OPE
coefficient at the boundary stratum~$D_{12}$; the logarithmic form
$\eta_{12}$ serves as a test form for the boundary pairing, not as a
multiplicative factor in a meromorphic expression.

\begin{remark}[Where the OPE data goes]\label{rem:frame-ope-extraction}
The map $d_{\mathrm{res}}\colon \bar{B}^{\mathrm{ch}}_1 \to
\bar{B}^{\mathrm{ch}}_0 = \mathbb{C}$ sends each pair of currents
to the constant $k$.  On $\mathbb{P}^1$, where $H^1(X) = 0$, the
kernel of $d_{\mathrm{res}}$ on $\bar{B}_1$ is trivial: every element
maps to a nonzero multiple of the vacuum.  At higher genus, zero modes
of $\alpha(z)$ along cycles of $\Sigma_g$ contribute nontrivially to
$\ker(d_{\mathrm{res}})$, producing $H^1 \cong H^1(\Sigma_g, \mathbb{C})$
(Theorem~\ref{thm:heisenberg-higher-genus}).
\end{remark}

% ================================================================
\section{The bar complex at degree 2 --- Arnold enters}
\label{sec:frame-bar-deg2}
% ================================================================

The degree-$2$ bar complex $\bar{B}^{\mathrm{ch}}_2$ involves three
currents and two logarithmic forms:
\begin{equation}\label{eq:frame-bar2-element}
\alpha(z_1) \otimes \alpha(z_2) \otimes \alpha(z_3)
\otimes \eta_{12} \wedge \eta_{23}.
\end{equation}

\subsubsection*{Sign conventions for the bar differential at degree~$2$}

Before computing, we fix signs.  The bar differential is a sum over
boundary strata:
\begin{equation}\label{eq:frame-d-sum-strata}
d_{\mathrm{res}} = \sum_{1 \leq i < j \leq n+1}
  (-1)^{\sigma(i,j)}\, \partial_{D_{ij}}
\end{equation}
where $\partial_{D_{ij}}$ restricts to the divisor
$D_{ij} = \{z_i = z_j\}$ and extracts the OPE coefficient, and
$\sigma(i,j)$ is the Koszul sign arising from permuting the
desuspended elements past each other.

For degree-$2$ elements with three desuspended currents
$s^{-1}\alpha(z_1) \otimes s^{-1}\alpha(z_2) \otimes s^{-1}\alpha(z_3)$,
the Koszul sign rule gives the following.
Each desuspended current $s^{-1}\alpha$ has degree~$2$
(conformal weight $1$ plus desuspension shift $1$), hence
is \emph{even} in the graded sense.  Moving even-degree elements
past each other produces sign $(-1)^{2 \cdot 2} = +1$.
Consequently, for the Heisenberg algebra (single even generator),
\emph{all Koszul signs from the coalgebra coproduct are~$+1$}.
The signs in the bar differential come entirely from the
orientation of the logarithmic forms on the boundary strata.

\subsubsection*{Term-by-term computation of $d_{\mathrm{res}}$}

The element~\eqref{eq:frame-bar2-element} lives on
$\overline{C}_3(X)$, whose boundary divisors are $D_{12}$, $D_{23}$,
and $D_{13}$.  The bar differential restricts to each:

\medskip
\noindent\textbf{Boundary stratum $D_{12}$: collision $z_1 = z_2$.}
The OPE $\alpha(z_1)\alpha(z_2) = k/(z_1 - z_2)^2 + \mathrm{reg.}$
collapses the pair to the scalar~$k$.  The form
$\eta_{12} \wedge \eta_{23}$ restricts to $D_{12}$ by the rule:
$\eta_{12}$ is ``consumed'' by the boundary pairing
on~$D_{12}$ (it is the test form for this stratum), and the
residual form is $\eta_{23}|_{z_1 = z_2}$.  Since $z_1 = z_2$
on $D_{12}$, the residual form is simply $\eta_{23}$ (which
does not involve $z_1$, so it is unaffected by the collision).
The sign from the wedge product is $+1$ because $\eta_{12}$ is
the \emph{left} factor:
\begin{equation}\label{eq:frame-D12-contribution}
\partial_{D_{12}}\bigl(
  \alpha \otimes \alpha \otimes \alpha \otimes
  \eta_{12} \wedge \eta_{23}
\bigr)
= +k \cdot \alpha(z_3) \otimes \eta_{23}.
\end{equation}
Here $\alpha(z_3) \otimes \eta_{23}$ is a degree-$1$
element: a single surviving current with a single logarithmic
form, living on $\overline{C}_2(X)$ with coordinates $(z_{12}, z_3)$
where $z_{12}$ is the collision point.

\medskip
\noindent\textbf{Boundary stratum $D_{23}$: collision $z_2 = z_3$.}
The OPE $\alpha(z_2)\alpha(z_3) = k/(z_2 - z_3)^2 + \mathrm{reg.}$
collapses the pair to~$k$.  The form $\eta_{12} \wedge \eta_{23}$
restricts to $D_{23}$ by consuming $\eta_{23}$ (the test form
for this stratum).  To consume $\eta_{23}$, we must move it to
the left of $\eta_{12}$, producing a sign from the wedge:
\[
\eta_{12} \wedge \eta_{23} = -\eta_{23} \wedge \eta_{12}.
\]
The residual form after consuming $\eta_{23}$ is $\eta_{12}|_{z_2 = z_3}$.
On $D_{23}$, $z_2 = z_3$, so $\eta_{12}|_{z_2 = z_3} = \eta_{13}$
(the form $d\log(z_1 - z_2)$ becomes $d\log(z_1 - z_3)$
when $z_2 = z_3$).  The surviving current is $\alpha(z_1)$.
Including the sign from the wedge permutation:
\begin{equation}\label{eq:frame-D23-contribution}
\partial_{D_{23}}\bigl(
  \alpha \otimes \alpha \otimes \alpha \otimes
  \eta_{12} \wedge \eta_{23}
\bigr)
= -k \cdot \alpha(z_1) \otimes \eta_{13}.
\end{equation}

\medskip
\noindent\textbf{Boundary stratum $D_{13}$: collision $z_1 = z_3$.}
The OPE $\alpha(z_1)\alpha(z_3) = k/(z_1 - z_3)^2 + \mathrm{reg.}$
collapses the pair to~$k$.  However, the $2$-form
$\eta_{12} \wedge \eta_{23}$ does \emph{not} contain the
factor $\eta_{13}$ as a wedge component.  The restriction of
$\eta_{12} \wedge \eta_{23}$ to $D_{13}$ involves:
\[
\eta_{12}\big|_{z_1 = z_3} = \eta_{32}
= -\eta_{23}, \qquad
\eta_{23}\big|_{z_1 = z_3} = \eta_{23}.
\]
So $(\eta_{12} \wedge \eta_{23})|_{D_{13}}
= (-\eta_{23}) \wedge \eta_{23} = 0$.
The boundary stratum $D_{13}$ contributes \emph{zero}:
\begin{equation}\label{eq:frame-D13-contribution}
\partial_{D_{13}}\bigl(
  \alpha \otimes \alpha \otimes \alpha \otimes
  \eta_{12} \wedge \eta_{23}
\bigr)
= 0.
\end{equation}
This vanishing is specific to the chosen basis element
$\eta_{12} \wedge \eta_{23}$; for a general $2$-form,
$D_{13}$ can contribute.

\medskip
\noindent\textbf{Assembling the differential.}
Summing the three contributions
\eqref{eq:frame-D12-contribution}--\eqref{eq:frame-D13-contribution}:
\begin{equation}\label{eq:frame-dres-deg2}
d_{\mathrm{res}}\bigl(\alpha(z_1) \otimes \alpha(z_2)
\otimes \alpha(z_3) \otimes \eta_{12} \wedge \eta_{23}\bigr)
= k \cdot \alpha(z_3) \otimes \eta_{23}
  - k \cdot \alpha(z_1) \otimes \eta_{13}.
\end{equation}

\subsubsection*{Verification: $d_{\mathrm{res}}^2 = 0$ at degree~$2$}

Now apply $d_{\mathrm{res}}$ a second time.  Each term in the
right-hand side of~\eqref{eq:frame-dres-deg2} is a degree-$1$
element, and $d_{\mathrm{res}}$ on degree~$1$ extracts $k$
by~\eqref{eq:frame-dres-deg1}:
\begin{equation}\label{eq:frame-dsquared-deg2}
d_{\mathrm{res}}^2 = k \cdot d_{\mathrm{res}}(\alpha(z_3)
\otimes \eta_{23}) - k \cdot d_{\mathrm{res}}(\alpha(z_1)
\otimes \eta_{13}) = k^2 - k^2 = 0.
\end{equation}
The first term: $d_{\mathrm{res}}(\alpha(z_3) \otimes \eta_{23})$
extracts the OPE of $\alpha$ with itself at the collision
$z_2 = z_3$, giving $k$.  The second term:
$d_{\mathrm{res}}(\alpha(z_1) \otimes \eta_{13})$ extracts the
OPE at $z_1 = z_3$, also giving $k$.  These are equal, and
the relative sign~$-1$ forces exact cancellation.

\subsubsection*{Alternative basis and the role of $D_{13}$}

To see that the vanishing of $D_{13}$ above was a feature
of the basis choice, consider the alternative degree-$2$ element
$\alpha^{\otimes 3} \otimes \eta_{23} \wedge \eta_{31}$.
The Arnold relation~\eqref{eq:frame-arnold} below gives
\[
\eta_{23} \wedge \eta_{31}
= -\eta_{12} \wedge \eta_{23} - \eta_{31} \wedge \eta_{12}
= -\eta_{12} \wedge \eta_{23} + \eta_{12} \wedge \eta_{31}.
\]
Computing $d_{\mathrm{res}}$ on this element, the stratum $D_{13}$
\emph{does} contribute (with coefficient $+k$), but the stratum
$D_{12}$ picks up an additional sign.  The computation of
$d_{\mathrm{res}}^2 = 0$ on \emph{any} basis element
requires the \emph{full} Arnold relation, not just pairwise
cancellation.  This is the essential content: the Arnold relation
is the universal identity guaranteeing $d^2 = 0$.

The cancellation is not an accident.  It is an instance of a general
identity on the logarithmic forms $\eta_{ij}$.

\begin{proposition}[Arnold relation]\label{prop:frame-arnold}
For any three distinct points $z_1, z_2, z_3 \in X$, the
logarithmic forms satisfy:
\begin{equation}\label{eq:frame-arnold}
\eta_{12} \wedge \eta_{23} + \eta_{23} \wedge \eta_{31}
+ \eta_{31} \wedge \eta_{12} = 0
\end{equation}
in $\Omega^2_{\log}(\overline{C}_3(X))$.
\end{proposition}

\begin{proof}
The identity $\frac{1}{z_1-z_2}\cdot\frac{1}{z_2-z_3}
+ \frac{1}{z_2-z_3}\cdot\frac{1}{z_3-z_1}
+ \frac{1}{z_3-z_1}\cdot\frac{1}{z_1-z_2} = 0$
follows from partial fractions: the three terms have
total numerator $(z_3-z_1)+(z_1-z_2)+(z_2-z_3) = 0$.
Taking $d\log$ of each factor and wedging produces
\eqref{eq:frame-arnold}.
\end{proof}

The Arnold relation forces $d^2 = 0$ on the bar complex.  At degree~$2$,
the two paths from $\bar{B}_2$ to $\bar{B}_0$ through $\bar{B}_1$
correspond to the two ``boundary of boundary'' strata in
$\overline{C}_3(X)$: $D_{12} \cap D_{23}$ and $D_{13} \cap D_{23}$.
The Arnold relation ensures these two paths contribute with opposite
signs, so they cancel.

More precisely: the differential $d$ on the bar complex is a sum of
residue maps along all collision divisors.  When $d$ is applied twice,
each pair of successive collisions defines a codimension-$2$ stratum
of $\overline{C}_{n+1}(X)$.  By Stokes' theorem on manifolds with
corners, $d^2 = 0$ if and only if the boundary of the boundary
vanishes, which is precisely the Arnold relation on the logarithmic
forms.

\begin{remark}[The Arnold relation as factorization coherence]
\label{rem:frame-arnold-factorization}
The identity~\eqref{eq:frame-arnold} is the fundamental coherence
relation of factorization algebras on algebraic curves.  It governs
how the chiral algebra structure extends from pairwise OPE data
(degree~$1$) to triple collisions (degree~$2$) and beyond.  Every
construction in this monograph rests on this identity: the bar complex
is well-defined because Arnold holds; Verdier duality exchanges bar
and cobar because Arnold is self-dual; the genus tower extends the
bar complex over $\overline{\mathcal{M}}_g$ because Arnold lifts to
the boundary strata of the Fulton--MacPherson operad.

In the language of operads, the Arnold relation is the quadratic
relation defining the \emph{gravity operad} $\mathrm{Grav}$
(the Koszul dual of the hypercommutative operad $\mathrm{HyCom}$),
which governs genus-$0$ chiral algebra structure.  The genus-$0$
bar-cobar adjunction is operadic Koszul duality:
$\mathrm{HyCom} \leftrightarrow \mathrm{Grav}$.
\end{remark}

% ================================================================
\section{The bar complex at all degrees --- $d^2 = 0$}
\label{sec:frame-bar-all}
% ================================================================

The pattern continues.  At degree~$n$, the bar complex element is
\begin{equation}\label{eq:frame-barn-element}
\alpha(z_1) \otimes \cdots \otimes \alpha(z_{n+1})
\otimes \eta_{i_1 j_1} \wedge \cdots \wedge \eta_{i_n j_n}
\end{equation}
and the bar differential extracts OPE data at each collision divisor.
Because the Heisenberg OPE has \emph{no simple pole}, the bracket
component $d_{\mathrm{bracket}}$ of the bar differential vanishes
identically.  Only the curvature component
$d_{\mathrm{curvature}}$ (which extracts the double-pole coefficient)
is nonzero.

Before stating the general theorem, we carry out the degree-$3$
computation in full detail.  This is the first degree where the
cancellation mechanism for $d^2 = 0$ involves \emph{pairs} of
Arnold relations, revealing the pattern that persists at all
higher degrees.

\subsection{Explicit degree-$3$ computation ($n = 3$, four points)}
\label{subsec:frame-degree3}

A degree-$3$ bar element lives on $\overline{C}_4(X)$ with four
points $z_1, z_2, z_3, z_4$ and involves a $3$-form in the
logarithmic algebra.  A natural basis element is:
\begin{equation}\label{eq:frame-bar3-element}
\xi_3 = \alpha(z_1) \otimes \alpha(z_2) \otimes \alpha(z_3)
  \otimes \alpha(z_4)
  \otimes \eta_{12} \wedge \eta_{23} \wedge \eta_{34}
  \;\in\; \bar{B}^{\mathrm{ch}}_3.
\end{equation}
This is the ``chain'' basis element, where the logarithmic forms
link consecutive pairs $1$--$2$, $2$--$3$, $3$--$4$.

\subsubsection*{Boundary strata of $\overline{C}_4(X)$}

The boundary divisors of $\overline{C}_4(X)$ are $D_{ij}$ for
$1 \leq i < j \leq 4$: there are $\binom{4}{2} = 6$ of them.
For each divisor $D_{ij}$, the bar differential extracts the
OPE $\alpha(z_i)\alpha(z_j) = k/(z_i - z_j)^2 + \mathrm{reg.}$,
collapses two of the four currents to the scalar $k$, and restricts
the $3$-form to a $2$-form on $\overline{C}_3(X)$.

We compute the contribution of each stratum.

\medskip
\noindent\textbf{Stratum $D_{12}$: collision $z_1 = z_2$.}
The form $\eta_{12}$ is consumed (it is the leftmost factor and
is the test form for $D_{12}$).  The OPE gives $k$.  The residual
$2$-form is $(\eta_{23} \wedge \eta_{34})|_{z_1 = z_2}$.  Since
$\eta_{23}$ and $\eta_{34}$ do not involve $z_1$ in a way that
changes under $z_1 = z_2$, the residual form is
$\eta_{23} \wedge \eta_{34}$ (with $z_1$ replaced by the
collision point $z_{12}$, but the forms only involve $z_2, z_3, z_4$).
The surviving currents are $\alpha(z_3)$ and $\alpha(z_4)$
(plus the collision point, which carries the scalar~$k$).
The sign is $+1$ (leftmost form consumed):
\begin{equation}\label{eq:frame-D12-deg3}
\partial_{D_{12}}(\xi_3)
= +k \cdot \alpha(z_3) \otimes \alpha(z_4)
  \otimes \eta_{23} \wedge \eta_{34}.
\end{equation}
(Here and below, we relabel the collision point: $z_{12}$ is
identified with $z_2$ on the residual configuration space.)

\medskip
\noindent\textbf{Stratum $D_{23}$: collision $z_2 = z_3$.}
The form $\eta_{23}$ must be consumed.  It sits in the middle
of the wedge product $\eta_{12} \wedge \eta_{23} \wedge \eta_{34}$.
Moving $\eta_{23}$ to the left past $\eta_{12}$ costs one
transposition sign:
\[
\eta_{12} \wedge \eta_{23} \wedge \eta_{34}
= -\eta_{23} \wedge \eta_{12} \wedge \eta_{34}.
\]
After consuming $\eta_{23}$, the residual $2$-form is
$\eta_{12}|_{z_2 = z_3} \wedge \eta_{34}|_{z_2 = z_3}
= \eta_{13} \wedge \eta_{34}$
(since $z_2 = z_3$ sends $\eta_{12} \mapsto \eta_{13}$
and leaves $\eta_{34}$ unchanged).
The surviving currents are $\alpha(z_1)$ and $\alpha(z_4)$:
\begin{equation}\label{eq:frame-D23-deg3}
\partial_{D_{23}}(\xi_3)
= -k \cdot \alpha(z_1) \otimes \alpha(z_4)
  \otimes \eta_{13} \wedge \eta_{34}.
\end{equation}

\medskip
\noindent\textbf{Stratum $D_{34}$: collision $z_3 = z_4$.}
The form $\eta_{34}$ is consumed.  It is the rightmost factor,
so moving it to the front costs two transpositions:
\[
\eta_{12} \wedge \eta_{23} \wedge \eta_{34}
= +\eta_{34} \wedge \eta_{12} \wedge \eta_{23}
\]
(two transpositions: $(-1)^2 = +1$; alternatively, a cyclic
permutation of three $1$-forms has sign $(-1)^{3-1} = +1$).
After consuming $\eta_{34}$, the residual $2$-form is
$\eta_{12} \wedge \eta_{23}|_{z_3 = z_4}
= \eta_{12} \wedge \eta_{23}$
(since $z_3 = z_4$ does not affect $\eta_{12}$ or $\eta_{23}$,
and we relabel the collision point as $z_3$).
The surviving currents are $\alpha(z_1)$ and $\alpha(z_2)$:
\begin{equation}\label{eq:frame-D34-deg3}
\partial_{D_{34}}(\xi_3)
= +k \cdot \alpha(z_1) \otimes \alpha(z_2)
  \otimes \eta_{12} \wedge \eta_{23}.
\end{equation}

\medskip
\noindent\textbf{Stratum $D_{13}$: collision $z_1 = z_3$.}
The form $\eta_{13}$ does not appear explicitly in the wedge
$\eta_{12} \wedge \eta_{23} \wedge \eta_{34}$.  We must restrict
each factor to $D_{13}$:
\begin{align*}
\eta_{12}\big|_{z_1 = z_3} &= \eta_{32} = -\eta_{23}, \\
\eta_{23}\big|_{z_1 = z_3} &= \eta_{23}, \\
\eta_{34}\big|_{z_1 = z_3} &= \eta_{34}.
\end{align*}
The restricted $3$-form is
$(-\eta_{23}) \wedge \eta_{23} \wedge \eta_{34}
= 0$
since $\eta_{23} \wedge \eta_{23} = 0$.
Therefore:
\begin{equation}\label{eq:frame-D13-deg3}
\partial_{D_{13}}(\xi_3) = 0.
\end{equation}

\medskip
\noindent\textbf{Stratum $D_{24}$: collision $z_2 = z_4$.}
Restrict each factor to $D_{24}$:
\begin{align*}
\eta_{12}\big|_{z_2 = z_4} &= \eta_{14}, \\
\eta_{23}\big|_{z_2 = z_4} &= \eta_{43} = -\eta_{34}, \\
\eta_{34}\big|_{z_2 = z_4} &= \eta_{34}.
\end{align*}
The restricted $3$-form is
$\eta_{14} \wedge (-\eta_{34}) \wedge \eta_{34}
= -\eta_{14} \wedge \eta_{34} \wedge \eta_{34} = 0$
since $\eta_{34} \wedge \eta_{34} = 0$.
Therefore:
\begin{equation}\label{eq:frame-D24-deg3}
\partial_{D_{24}}(\xi_3) = 0.
\end{equation}

\medskip
\noindent\textbf{Stratum $D_{14}$: collision $z_1 = z_4$.}
Restrict each factor to $D_{14}$:
\begin{align*}
\eta_{12}\big|_{z_1 = z_4} &= \eta_{42} = -\eta_{24}, \\
\eta_{23}\big|_{z_1 = z_4} &= \eta_{23}, \\
\eta_{34}\big|_{z_1 = z_4} &= \eta_{34}.
\end{align*}
The restricted $3$-form is
$(-\eta_{24}) \wedge \eta_{23} \wedge \eta_{34}
= \eta_{24} \wedge \eta_{32} \wedge \eta_{34}$
after rewriting $-\eta_{23} = \eta_{32}$.
This does not manifestly vanish; however, for the boundary
pairing on $D_{14}$ to extract a nonzero contribution, the
$3$-form must contain $\eta_{14}$ as a factor.  Since $\eta_{14}$
does not appear (the form involves $\eta_{24}, \eta_{32}, \eta_{34}$
only), the boundary pairing gives zero.

Alternatively: after the collision $z_1 = z_4$, the form must
reduce to a $2$-form on $\overline{C}_3(X)$.  The restriction
$(-\eta_{24}) \wedge \eta_{23} \wedge \eta_{34}$ involves three
independent forms on a space of dimension $\leq 2$ (the log-form
space on three points is $2$-dimensional by Arnold), so it is a
relation, not a new element.  The pairing with the test form
$\eta_{14}$ extracts zero.
\begin{equation}\label{eq:frame-D14-deg3}
\partial_{D_{14}}(\xi_3) = 0.
\end{equation}

\subsubsection*{Assembling $d_{\mathrm{res}}$ at degree $3$}

Only the three ``adjacent'' strata $D_{12}$, $D_{23}$, $D_{34}$
contribute.  The three ``non-adjacent'' strata $D_{13}$, $D_{24}$,
$D_{14}$ all vanish by form-degree reasons.
Summing~\eqref{eq:frame-D12-deg3}--\eqref{eq:frame-D34-deg3}:
\begin{equation}\label{eq:frame-dres-deg3-full}
\boxed{\begin{aligned}
d_{\mathrm{res}}(\xi_3)
={} & k \cdot \alpha(z_3) \otimes \alpha(z_4)
    \otimes \eta_{23} \wedge \eta_{34} \\
  & {} - k \cdot \alpha(z_1) \otimes \alpha(z_4)
    \otimes \eta_{13} \wedge \eta_{34} \\
  & {} + k \cdot \alpha(z_1) \otimes \alpha(z_2)
    \otimes \eta_{12} \wedge \eta_{23}.
\end{aligned}}
\end{equation}
Observe the alternating signs $+, -, +$ from the position of the
consumed form in the wedge product: leftmost ($+1$), middle
($-1$), rightmost ($+1$).

\subsubsection*{Verification: $d_{\mathrm{res}}^2 = 0$ at degree~$3$}

We now apply $d_{\mathrm{res}}$ again to each of the three terms
in~\eqref{eq:frame-dres-deg3-full}.
Each term is a degree-$2$ element, and from
\eqref{eq:frame-dres-deg2} we know how $d_{\mathrm{res}}$ acts
at degree~$2$.

\medskip
\noindent\textbf{First term}: $d_{\mathrm{res}}\bigl(
k \cdot \alpha(z_3) \otimes \alpha(z_4)
\otimes \eta_{23} \wedge \eta_{34}\bigr)$.

Applying the degree-$2$ formula (with indices shifted
$1 \to 2, 2 \to 3, 3 \to 4$):
\begin{equation}\label{eq:frame-dsq-term1}
d_{\mathrm{res}}\bigl(\alpha(z_3) \otimes \alpha(z_4)
  \otimes \eta_{23} \wedge \eta_{34}\bigr)
= k \cdot \alpha(z_4) \otimes \eta_{34}
  - k \cdot \alpha(z_2) \otimes \eta_{24}.
\end{equation}
Wait --- this is not quite right.  The element
$\alpha(z_3) \otimes \alpha(z_4) \otimes \eta_{23} \wedge \eta_{34}$
lives on $\overline{C}_3$ with coordinates $(z_2, z_3, z_4)$
(the point $z_2$ appears through the form $\eta_{23}$ but carries
no current factor).  We must be careful: the bar differential
on this degree-$2$ element has boundary strata $D_{23}$
and $D_{34}$.

At $D_{23}$: the current at $z_3$ collides with the scalar at
$z_2$.  Since there is no current at $z_2$, this stratum extracts
no OPE --- it contributes zero.  More precisely: the degree-$2$
element has only two currents ($\alpha(z_3)$ and $\alpha(z_4)$),
so the only OPE extraction is at $D_{34}$:
\begin{equation}\label{eq:frame-dsq-term1-corrected}
d_{\mathrm{res}}\bigl(
  \alpha(z_3) \otimes \alpha(z_4)
  \otimes \eta_{23} \wedge \eta_{34}\bigr)
= k \cdot \eta_{23}|_{z_3 = z_4}
= k \cdot \eta_{24}.
\end{equation}
But $\eta_{23}$ is consumed by the form pairing, so we
also get a contribution from stratum $D_{23}$.  Let us
redo this carefully.

The element $k \cdot \alpha(z_3) \otimes \alpha(z_4)
\otimes \eta_{23} \wedge \eta_{34}$ is a degree-$2$ element
on $\overline{C}_3$ with coordinates $(z_2, z_3, z_4)$.
But it has only $2$ currents (at $z_3$ and $z_4$), not $3$.
This means it is not a generic degree-$2$ element:
it is a degree-$2$ element of the bar complex where one
slot carries the vacuum (scalar $k$ from the previous
contraction).  The bar differential on such an element
extracts the OPE of $\alpha(z_3)$ with $\alpha(z_4)$
at stratum $D_{34}$, consuming $\eta_{34}$.  The sign
from moving $\eta_{34}$ past $\eta_{23}$ is $-1$:
\begin{equation}\label{eq:frame-dsq-T1}
d_{\mathrm{res}}\bigl(
  k \cdot \alpha(z_3) \otimes \alpha(z_4)
  \otimes \eta_{23} \wedge \eta_{34}\bigr)
= k \cdot (-1) \cdot k \cdot \eta_{23}\big|_{z_3 = z_4}
= -k^2 \cdot \eta_{24}.
\end{equation}
After the collapse $z_3 = z_4$, no currents remain:
$-k^2 \cdot \eta_{24}$ is a $1$-form on $\overline{C}_2$
with coordinates $(z_2, z_{34})$, carrying no current factor
(a degree-$1$ ``scalar'' element).  But this is \emph{not}
a bar element in the usual sense; it is an element of
$\bar{B}_1^{\mathrm{ch}}$ tensored with a scalar coefficient.
The key point is simply its sign and coefficient.

\medskip
\noindent\textbf{Second term}: $d_{\mathrm{res}}\bigl(
- k \cdot \alpha(z_1) \otimes \alpha(z_4)
\otimes \eta_{13} \wedge \eta_{34}\bigr)$.

The only OPE-producing stratum is $D_{14}$ (collision of $z_1$
and $z_4$).  The form $\eta_{13} \wedge \eta_{34}$ must yield
$\eta_{14}$ --- but $\eta_{14}$ does not appear as a factor.
Restricting to $D_{14}$:
\[
\eta_{13}\big|_{z_1 = z_4} = \eta_{43} = -\eta_{34},
\qquad
\eta_{34}\big|_{z_1 = z_4} = \eta_{34}.
\]
So $(\eta_{13} \wedge \eta_{34})|_{D_{14}}
= (-\eta_{34}) \wedge \eta_{34} = 0$.
The contribution is zero.

There is also stratum $D_{34}$: the OPE is between
$\alpha(z_4)$ and... but there is no current at $z_3$
(only a form $\eta_{13}$ referencing $z_3$).  Since the
currents are at $z_1$ and $z_4$ only, the available collision
is $D_{14}$, which vanishes as shown.  Similarly,
$D_{13}$ has no current pair.

However, we must reconsider: this degree-$2$ element
has currents at $z_1$ and $z_4$, with forms $\eta_{13}$
and $\eta_{34}$ referencing $z_3$.  The bar differential
extracts the OPE of $\alpha(z_1)$ with $\alpha(z_4)$.
The relevant boundary stratum is $D_{14}$, and the test
form is $\eta_{14}$.  Since $\eta_{14}$ does not appear
as a factor, we use the Arnold relation:
\[
\eta_{13} \wedge \eta_{34}
= \eta_{13} \wedge \eta_{34}
\]
which, by the Arnold identity on points $(1, 3, 4)$, gives
$\eta_{13} \wedge \eta_{34}
= -\eta_{34} \wedge \eta_{41} - \eta_{41} \wedge \eta_{13}
= \eta_{34} \wedge \eta_{14} + \eta_{14} \wedge \eta_{13}$
(using $\eta_{41} = -\eta_{14}$).  So:
\[
\eta_{13} \wedge \eta_{34}
= \eta_{34} \wedge \eta_{14} + \eta_{14} \wedge \eta_{13}
= -\eta_{14} \wedge \eta_{34} + \eta_{14} \wedge \eta_{13}
= \eta_{14} \wedge (\eta_{13} - \eta_{34}).
\]
Now the factor $\eta_{14}$ appears, and the bar differential
extracts:
\begin{equation}\label{eq:frame-dsq-T2}
\begin{aligned}
d_{\mathrm{res}}\bigl(
  -k \cdot \alpha(z_1) \otimes \alpha(z_4)
  \otimes \eta_{13} \wedge \eta_{34}\bigr)
&= -k^2 (\eta_{13} - \eta_{34})\big|_{z_1 = z_4} \\
&= -k^2 (\eta_{43} - \eta_{34})
= -k^2(-\eta_{34} - \eta_{34})
= +2k^2 \eta_{34}.
\end{aligned}
\end{equation}
But wait: after the collapse $z_1 = z_4$, the residual
form should reference the collision point.  Let us write
$z_{14}$ for the collision point.  Then
$\eta_{13}|_{z_1 = z_4} = \eta_{(14),3}$ (distance from
collision to $z_3$) and
$\eta_{34}|_{z_1 = z_4} = \eta_{3,(14)}$ (same, with opposite
sign).  So $\eta_{13} - \eta_{34}$ becomes
$\eta_{(14),3} - \eta_{3,(14)} = 2\eta_{(14),3}$.

The key observation is that this intermediate computation is
not needed for verifying $d^2 = 0$.  What matters is that the
total of all three terms cancels.  We take a cleaner approach.

\subsubsection*{Clean verification of $d^2 = 0$ via Arnold pairs}

Rather than tracking the intermediate expressions through
relabeling, we verify $d^2 = 0$ structurally.  The bar differential
$d = d_{\mathrm{res}}$ satisfies $d^2 = 0$ if and only if
every codimension-$2$ boundary stratum of $\overline{C}_4(X)$
contributes zero.  The codimension-$2$ strata are of two types:

\medskip
\noindent\textbf{Type I: nested collisions} $D_{ij} \cap D_{jk}$
(three points $z_i, z_j, z_k$ collide in sequence).
These contribute via the composition
$\partial_{D_{jk}} \circ \partial_{D_{ij}}$:
first $z_i$ collides with $z_j$, then the collision point
collides with $z_k$.  The reverse order
$\partial_{D_{ij}} \circ \partial_{D_{jk}}$ gives the same
codimension-$2$ stratum with opposite orientation.  By the
Arnold relation on the triple $(z_i, z_j, z_k)$, these two
contributions cancel.

For the four-point configuration, the nested triples are:
$(1,2,3)$, $(2,3,4)$, $(1,2,4)$, $(1,3,4)$, $(2,3,4)$.
Each produces one Arnold-relation cancellation.

\medskip
\noindent\textbf{Type II: disjoint collisions} $D_{ij} \cap D_{kl}$
where $\{i,j\} \cap \{k,l\} = \emptyset$.
For four points, the only such stratum is $D_{12} \cap D_{34}$
(simultaneous collision $z_1 = z_2$ and $z_3 = z_4$).  The two
orderings $\partial_{D_{34}} \circ \partial_{D_{12}}$ and
$\partial_{D_{12}} \circ \partial_{D_{34}}$ commute (they act
on disjoint pairs), and the sign from the form restriction
ensures they cancel.  Explicitly, for the form
$\eta_{12} \wedge \eta_{23} \wedge \eta_{34}$:
\begin{itemize}
\item $\partial_{D_{12}}$ then $\partial_{D_{34}}$:
  consume $\eta_{12}$ (sign $+1$), leaving
  $\eta_{23} \wedge \eta_{34}$; then consume $\eta_{34}$
  (sign $-1$ from moving past $\eta_{23}$).  Net: $(-1)$.
\item $\partial_{D_{34}}$ then $\partial_{D_{12}}$:
  consume $\eta_{34}$ (sign $+1$ by cyclic permutation),
  leaving $\eta_{12} \wedge \eta_{23}$; then consume
  $\eta_{12}$ (sign $+1$).  Net: $(+1)$.
\end{itemize}
The two orderings produce opposite signs:
\begin{equation}\label{eq:frame-disjoint-cancel}
\partial_{D_{34}} \circ \partial_{D_{12}}(\xi_3)
+ \partial_{D_{12}} \circ \partial_{D_{34}}(\xi_3)
= k^2 \cdot (-1 + 1) = 0.
\end{equation}

\medskip
Since every codimension-$2$ stratum contributes zero (Type~I
by Arnold, Type~II by sign cancellation), we conclude:
\begin{equation}\label{eq:frame-dsquared-deg3}
\boxed{d_{\mathrm{res}}^2(\xi_3) = 0.}
\end{equation}

\subsubsection*{The adjacency pattern}

The degree-$3$ computation reveals a structural feature.  For
the chain basis element
$\alpha^{\otimes 4} \otimes \eta_{12} \wedge \eta_{23} \wedge \eta_{34}$,
only ``adjacent'' strata (those $D_{ij}$ where $\eta_{ij}$ appears
as a factor in the wedge) give nonzero contributions to~$d$.
The non-adjacent strata $D_{13}$, $D_{24}$, $D_{14}$ all vanish
by form-squaring ($\eta \wedge \eta = 0$).  This adjacency
pattern persists at all degrees for the chain basis: the
differential on
$\alpha^{\otimes(n+1)} \otimes \eta_{12} \wedge \eta_{23}
\wedge \cdots \wedge \eta_{n,n+1}$
receives contributions only from $D_{12}, D_{23}, \ldots, D_{n,n+1}$,
producing an alternating sum of $n$ terms with signs
$+1, -1, +1, -1, \ldots$

For a general basis element (not the chain basis), all strata can
contribute, and $d^2 = 0$ requires the full Arnold relation.
The chain basis is special precisely because it diagonalizes the
``adjacency'' of the logarithmic forms.

\subsection{The general theorem}

\begin{theorem}[Heisenberg bar complex at genus~$0$]
\label{thm:frame-heisenberg-bar}
For $\mathcal{H}_k$ on $\mathbb{P}^1$:
\begin{equation}\label{eq:frame-heisenberg-bar-cohomology}
H^n(\bar{B}^{\mathrm{ch}}_{\mathrm{geom}}(\mathcal{H}_k))
= \begin{cases}
\mathbb{C} & n = 0, \\
0 & n = 1, \\
\mathbb{C} \cdot c_k & n = 2, \\
0 & n > 2.
\end{cases}
\end{equation}
The class $c_k \in H^2$ arises from the triple-collision residue
and is the homological manifestation of the central extension.
The differential satisfies $d^2 = 0$ at all degrees by the
Arnold relation.
\end{theorem}

\begin{proof}
\emph{Degree~$0$.}
$\bar{B}^0 = \mathbb{C} \cdot 1$ (vacuum), $d = 0$.

\emph{Degree~$1$.}
The differential $d_{\mathrm{res}}\colon \bar{B}^1 \to \bar{B}^0$
sends $\alpha(z_1) \otimes \alpha(z_2) \otimes \eta_{12} \mapsto k \cdot 1$
(equation~\eqref{eq:frame-dres-deg1}).
On $\mathbb{P}^1$, where $H^1(X) = 0$, this map is surjective
and has trivial kernel, so $H^1 = 0$.

\emph{Degree~$2$.}
The three-current collision produces the central charge class.
When $\alpha(z_1), \alpha(z_2), \alpha(z_3)$ approach a
common point, the iterated OPE produces:
\begin{equation}\label{eq:frame-triple-residue}
\mathrm{Res}_{D_{123}}\bigl[\alpha(z_1)\alpha(z_2)\alpha(z_3)
\otimes \eta_{12} \wedge \eta_{23}\bigr]
= k \cdot [\text{central class}]
\end{equation}
The commutator $[\alpha(z), \alpha(w)] = k \cdot
\partial_w\delta(z-w)$ contributes a central term at
the triple collision.  This produces a nontrivial
class $c_k \in H^2$.

\emph{Degrees $\geq 3$.}
Vanish by dimension counting: the Heisenberg algebra is
generated by a single current of conformal weight~$1$, and the
bar complex cohomology is controlled by the Lie algebra cohomology
of the abelian Lie algebra $\mathfrak{h}$, which is concentrated
in degrees $\leq 2$ for a $1$-dimensional central extension.

\emph{$d^2 = 0$.}
At each degree, $d^2 = 0$ by the Arnold relation
(Proposition~\ref{prop:frame-arnold}) applied to the
logarithmic forms on $\overline{C}_{n+1}(X)$.
The degree-$2$ verification is~\eqref{eq:frame-dsquared-deg2};
the general case follows by the same argument on all codimension-$2$
strata of $\overline{C}_{n+1}(X)$.
\end{proof}

The explicit degree-$2$ and degree-$3$ bar differentials confirm the
pattern.  The degree-$2$ computation is~\eqref{eq:frame-dres-deg2}
with $d^2 = 0$ by~\eqref{eq:frame-dsquared-deg2}; the degree-$3$
computation is~\eqref{eq:frame-dres-deg3-full} with $d^2 = 0$
by~\eqref{eq:frame-dsquared-deg3}.
At degree~$3$, the cancellation mechanism for $d^2 = 0$
involves both nested Arnold relations (Type~I) and disjoint-collision
sign cancellations (Type~II), as detailed in
\S\ref{subsec:frame-degree3}.

% ================================================================
\section{Reading off the Koszul dual}\label{sec:frame-koszul-dual}
% ================================================================

The bar complex of $\mathcal{H}_k$ is a coalgebra.  Its coproduct
structure is determined by the OPE.  Because the double-pole
residue $\mathrm{Res}_{z_i = z_j}[\omega/(z_i - z_j)^2]
= \partial_{z_i}\omega|_{z_i = z_j}$ is symmetric under
$z_i \leftrightarrow z_j$, the coproduct is \emph{cocommutative}:
\begin{equation}\label{eq:frame-primitive}
\Delta(\alpha) = \alpha \otimes 1 + 1 \otimes \alpha.
\end{equation}
By the Milnor--Moore theorem, a conilpotent cocommutative coalgebra
is a coLie coalgebra.  The bar complex is therefore:
\begin{equation}\label{eq:frame-bar-coalgebra}
\bar{B}^{\mathrm{ch}}(\mathcal{H}_k)
\simeq \mathrm{coLie}^{\mathrm{ch}}(V^*)
\end{equation}
the conilpotent chiral Lie coalgebra on the single cogenerator
$V^* = \mathrm{Span}\{\alpha^*\}$.

Apply the cobar functor.  The cobar of a coLie coalgebra is the
free commutative algebra (this is the dual statement to
$\mathrm{Com}^! = \mathrm{Lie}$ in the Koszul duality of operads):
\begin{equation}\label{eq:frame-cobar}
\Omega^{\mathrm{ch}}\bigl(\mathrm{coLie}^{\mathrm{ch}}(V^*)\bigr)
\simeq \mathrm{Sym}^{\mathrm{ch}}(V^*).
\end{equation}

The level $k$ determines a curvature element
$m_0 = -k \cdot \omega \in \mathrm{Sym}^2(V^*)$, encoding the
central extension.

\begin{theorem}[Heisenberg Koszul dual]
\label{thm:frame-heisenberg-koszul-dual}
The Koszul dual of $\mathcal{H}_k$ is the curved commutative
chiral algebra
\begin{equation}\label{eq:frame-koszul-dual}
\mathcal{H}_k^! \simeq \bigl(\mathrm{Sym}^{\mathrm{ch}}(V^*),\;
d = 0,\; m_0 = -k \cdot \omega\bigr).
\end{equation}
This is \emph{not} $\mathcal{H}_{-k}$: the Heisenberg algebra is
not self-dual.
\end{theorem}

The distinction deserves emphasis:
\begin{center}
\small
\begin{tabular}{c|c}
\textbf{Heisenberg} $\mathcal{H}_k$ &
\textbf{Koszul dual} $\mathrm{Sym}^{\mathrm{ch}}(V^*)$ \\
\hline
Non-commutative & Commutative \\
$[\alpha, \alpha] = k \neq 0$ &
$\alpha \cdot \alpha = \alpha^2$ (commutes) \\
Central extension & No central extension \\
Double-pole OPE & Regular OPE \\
Lie-type & Com-type
\end{tabular}
\end{center}

The construction we have just performed for $\mathcal{H}_k$
is an instance of a general pattern: the bar complex of any
augmented chiral algebra is a coalgebra, and the cobar of
that coalgebra produces the Koszul dual.  The general
bar-cobar adjunction is developed in
Chapter~\ref{chap:bar-cobar}.

\begin{remark}[Three different ``dualities'']
\label{rem:frame-three-dualities}
The literature contains three distinct duality operations
on the Heisenberg algebra, which must not be confused:
\begin{enumerate}
\item \emph{Bar-cobar Koszul duality} (this chapter):
  $\mathcal{H}_k \leftrightarrow \mathrm{Sym}^{\mathrm{ch}}(V^*)$.
  Exchanges Lie-type and Com-type structure.
\item \emph{Level-shifting duality}:
  $\mathrm{Rep}(\mathcal{H}_k)
  \simeq \mathrm{Rep}(\mathcal{H}_{-k})$.
  An equivalence of representation categories.
\item \emph{Boson-fermion correspondence}:
  $\mathrm{Rep}(\mathcal{H}_k)
  \simeq \mathrm{Rep}(\mathcal{F}^{\otimes 2})$.
  Relates Heisenberg modules to free-fermion modules.
\end{enumerate}
Only~(1) is the subject of this monograph.
\end{remark}

% ================================================================
\section{Bar-cobar inversion --- Theorem A in action}
\label{sec:frame-inversion}
% ================================================================

The cobar construction recovers the original algebra:
\begin{equation}\label{eq:frame-bar-cobar-inversion}
\Omega^{\mathrm{ch}}\bigl(\bar{B}^{\mathrm{ch}}(\mathcal{H}_k)\bigr)
\xrightarrow{\;\sim\;} \mathcal{H}_k.
\end{equation}

The chain of identifications is:
\[
\bar{B}^{\mathrm{ch}}(\mathcal{H}_k)
\simeq \mathrm{coLie}^{\mathrm{ch}}(V^*)
\quad\Longrightarrow\quad
\Omega^{\mathrm{ch}}(\mathrm{coLie}^{\mathrm{ch}}(V^*))
\simeq \mathrm{Sym}^{\mathrm{ch}}(V^*)
\]
and applying bar again:
\[
\bar{B}^{\mathrm{ch}}(\mathrm{Sym}^{\mathrm{ch}}(V^*))
\simeq \mathrm{coLie}^{\mathrm{ch}}(V)
\xrightarrow{\;\Omega\;}
\mathcal{H}_k.
\]

The last step uses the fact that the free Lie algebra on a single
generator is $1$-dimensional (abelian), so the cobar of the coLie
coalgebra on one cogenerator recovers the original Heisenberg
algebra.  The bar-cobar adjunction $\Omega \dashv \bar{B}$ restricts
to an equivalence on Koszul algebras; $\mathcal{H}_k$ is Koszul
precisely because its bar complex is formal (concentrated in
degrees $0$ and $2$).

This is Theorem~A (Theorem~\ref{thm:bar-cobar-isomorphism-main})
for the simplest non-trivial case.  The general theorem establishes
bar-cobar duality for all augmented chiral algebras, with the
quasi-isomorphism $\Omega \bar{B}(\mathcal{A}) \xrightarrow{\sim}
\mathcal{A}$ holding on the Koszul locus.  We will see in
Chapter~\ref{chap:bar-cobar} that this locus includes
all affine Kac--Moody algebras and $\mathcal{W}$-algebras whose
PBW spectral sequence collapses.

% ================================================================
\section{The passage to genus~1 --- curvature appears}
\label{sec:frame-genus1}
% ================================================================

Everything computed so far takes place on $\mathbb{P}^1$
(genus~$0$).  Replace $X = \mathbb{P}^1$ by an elliptic curve
$E_\tau = \mathbb{C}/(\mathbb{Z} + \tau\mathbb{Z})$ with modular
parameter $\tau \in \mathbb{H}$.

The propagator changes.  On $\mathbb{P}^1$, the logarithmic form
$\eta_{12} = d\log(z_1 - z_2)$ has a simple pole at $z_1 = z_2$
and is single-valued.  On $E_\tau$, no meromorphic function with
a single simple pole can be doubly periodic (Liouville's theorem).
The correct replacement is the \emph{elliptic propagator}:
\begin{equation}\label{eq:frame-elliptic-propagator}
G_\tau(z) = \zeta_\tau(z) + \frac{\pi^2 E_2(\tau)}{3}\, z
\end{equation}
where $\zeta_\tau(z)$ is the Weierstrass zeta function
\begin{equation}\label{eq:frame-weierstrass-zeta}
\zeta_\tau(z) = \frac{1}{z}
+ \sum_{(m,n) \neq (0,0)}\left[
  \frac{1}{z - m - n\tau}
  + \frac{1}{m + n\tau}
  + \frac{z}{(m+n\tau)^2}
\right]
\end{equation}
and $E_2(\tau)$ is the weight-$2$ Eisenstein series
\begin{equation}\label{eq:frame-E2}
E_2(\tau) = 1 - 24\sum_{n=1}^{\infty} \sigma_1(n)\, q^n,
\qquad q = e^{2\pi i \tau}.
\end{equation}

The Weierstrass $\zeta$-function is \emph{quasi-periodic}: it
satisfies $\zeta_\tau(z + 1) = \zeta_\tau(z) + 2\eta_1$ and
$\zeta_\tau(z + \tau) = \zeta_\tau(z) + 2\eta_2$, where the
quasi-periods are $\eta_1 = \pi^2 E_2(\tau)/6$.  The
$E_2$-correction in~\eqref{eq:frame-elliptic-propagator}
is the price of modular covariance: without it, the propagator
does not transform correctly under $\mathrm{SL}_2(\mathbb{Z})$.

The Laurent expansion near $z = 0$ is:
\begin{equation}\label{eq:frame-green-laurent}
G_\tau(z) = \frac{1}{z}
+ \frac{\pi^2 E_2(\tau)}{3}\, z + O(z^3).
\end{equation}
At genus~$0$, only the pole $1/z$ survives.  The $E_2$-term
is the first appearance of modular forms in the bar complex.

\subsection{The genus-1 two-point function}

The Heisenberg two-point function on $E_\tau$ is
\begin{equation}\label{eq:frame-genus1-2pt}
\langle \alpha(z_1)\, \alpha(z_2) \rangle_{E_\tau}
= k \cdot G_\tau(z_1 - z_2).
\end{equation}

The genus-$1$ partition function is
\begin{equation}\label{eq:frame-partition-g1}
Z_{E_\tau}^{\mathcal{H}} =
\mathrm{Tr}_{\mathcal{F}_k}\, q^{L_0 - c/24}
= \frac{1}{\eta(\tau)}
= q^{-1/24}\prod_{n=1}^{\infty}\frac{1}{1-q^n}
\end{equation}
where $\eta(\tau)$ is the Dedekind eta function.  The coefficients
are the partition numbers $p(n)$: the number of partitions of $n$
counts the number of states at excitation level~$n$ in the
Fock space.

\subsection{Curvature: $d^2 \neq 0$ at genus 1}

Compute the bar differential on $E_\tau$ using the elliptic
propagator~\eqref{eq:frame-elliptic-propagator}.  The B-cycle
monodromy of $G_\tau$ produces a defect: transporting the
propagator around the B-cycle of $E_\tau$ shifts it by the
constant $-2\pi i$.  This shift is the genus-$1$ curvature.

\begin{theorem}[Genus-1 curvature]\label{thm:frame-genus1-curvature}
For $\mathcal{H}_k$ on $E_\tau$, the bar differential satisfies
\begin{equation}\label{eq:frame-dsquared-genus1}
d^2 = k \cdot \omega_1 \cdot \mathrm{id}
\end{equation}
where $\omega_1 = c_1(\mathbb{E}) \in H^2(\overline{\mathcal{M}}_1)$
is the first Chern class of the Hodge bundle on the moduli space
of elliptic curves.
\end{theorem}

The proof is the direct computation: the B-cycle monodromy of the
elliptic propagator contributes a phase $e^{2\pi i k}$ to the
parallel transport of bar complex elements.  Squaring the
differential, the two B-cycle transports contribute a factor
$(-2\pi i)^2$ that, after integration over $\overline{\mathcal{M}}_1$,
produces the Hodge class $\omega_1 = c_1(\mathbb{E})$ with
coefficient~$k$.  The detailed computation appears in
Theorem~\ref{thm:heisenberg-obs}.

The differential has ceased to be strictly square-zero.

At genus~$0$, the bar complex was a genuine chain complex
($d^2 = 0$) and its cohomology was well-defined.  At genus~$1$,
the bar complex is \emph{curved}: $d^2 = k \cdot \omega_1$
is proportional to the level.  The curvature vanishes only at
$k = 0$, where the Heisenberg algebra degenerates.

\begin{remark}[Physical meaning of curvature]
\label{rem:frame-curvature-physics}
In conformal field theory, the curvature $d^2 = k \cdot \omega_1$
is the \emph{conformal anomaly} at one loop.  The level $k$
controls the strength of the central extension; the Hodge class
$\omega_1$ is the gravitational coupling.  The product $k \cdot
\omega_1$ measures the obstruction to defining the chiral algebra
consistently on all elliptic curves simultaneously: it is the
first Chern class of the \emph{determinant line bundle}
$\det(R\pi_*\mathcal{F}_k)$ on $\overline{\mathcal{M}}_1$,
where $\pi\colon \mathcal{E} \to \overline{\mathcal{M}}_1$
is the universal elliptic curve and $\mathcal{F}_k$ is the
Fock space sheaf.
\end{remark}

% ================================================================
\section{The holomorphic anomaly and $E_2$}
\label{sec:frame-holomorphic-anomaly}
% ================================================================

The appearance of $E_2$ introduces a \emph{holomorphic anomaly}.
Under modular transformations $\tau \mapsto (a\tau+b)/(c\tau+d)$,
the weight-$4$ and weight-$6$ Eisenstein series transform as
modular forms:
\[
E_4(\gamma \cdot \tau) = (c\tau+d)^4 E_4(\tau), \qquad
E_6(\gamma \cdot \tau) = (c\tau+d)^6 E_6(\tau).
\]
The weight-$2$ Eisenstein series does not:
\begin{equation}\label{eq:frame-E2-anomaly}
E_2\!\left(\frac{a\tau+b}{c\tau+d}\right)
= (c\tau+d)^2 E_2(\tau) + \frac{6c(c\tau+d)}{\pi i}.
\end{equation}
The extra term $6c(c\tau+d)/(\pi i)$ is the holomorphic anomaly.

For the Heisenberg two-point function on $E_\tau$, this anomaly
produces:
\begin{equation}\label{eq:frame-2pt-anomaly}
\langle \alpha(z_1)\, \alpha(z_2) \rangle_{\gamma \cdot \tau}
= (c\tau+d)^2
  \langle \alpha((c\tau+d)^{-1}z_1)\,
  \alpha((c\tau+d)^{-1}z_2) \rangle_\tau
  - \frac{6ck}{\pi i}\, z_{12}.
\end{equation}
The anomaly term is proportional to $k$ (the level) and to $c$
(the lower-left entry of the modular matrix).

This is precisely the genus-$1$ obstruction class
(Theorem~\ref{thm:frame-genus1-curvature}), seen from the modular
side.  The bar complex curvature $d^2 = k \cdot \omega_1$ and the
holomorphic anomaly in $E_2$ are two manifestations of the same
phenomenon: the failure of the Heisenberg algebra to extend
coherently over the moduli space $\overline{\mathcal{M}}_1$.

% ================================================================
\section{Why the ordinary derived category fails}
\label{sec:frame-coderived}
% ================================================================

The curvature $d^2 = k \cdot \omega_1$ means that the genus-$1$
bar complex is not a chain complex in the ordinary sense: it does
not live in the derived category $D(\mathrm{Vect})$.  It lives in
the \emph{coderived category} $D^{\mathrm{co}}(\mathrm{Vect})$ of
Positselski~\cite{Positselski11}.

The distinction is subtle but essential.  In the ordinary derived
category, a complex with $d^2 \neq 0$ is simply undefined.  In
Positselski's framework, a \emph{curved} differential graded
(co)algebra --- one where $d^2 = m_0 \neq 0$ for some curvature
element $m_0$ --- has a well-defined coderived category in which:
\begin{enumerate}
\item The bar-cobar adjunction still holds (as an adjunction of
  curved (co)algebras).
\item The notion of ``acyclic'' is refined: a curved complex can
  be \emph{coacyclic} (contractible in the coderived sense) without
  being acyclic in the ordinary sense.
\item The bar-cobar quasi-isomorphism
  $\Omega(\bar{B}(\mathcal{A})) \to \mathcal{A}$ holds in the
  coderived category when it fails in the ordinary derived category.
\end{enumerate}

The Heisenberg algebra at $k \neq 0$ is the simplest object that
forces us out of the ordinary derived category.  At $k = 0$,
$d^2 = 0$ and the ordinary derived category suffices; at $k \neq 0$,
the curvature is nonzero and only the coderived category gives the
correct homological algebra.

This is not an abstract concern.  The genus-$1$ bar complex of
$\mathcal{H}_k$ with $k \neq 0$ is ``acyclic'' in the na\"ive sense
(its homology with respect to $d$ is zero, since $d^2 \neq 0$
means the notion of ``homology'' is undefined), but it is
\emph{not coacyclic} in Positselski's sense --- it carries
genuine invariants (the Hodge class $\omega_1$, the partition
function $1/\eta(\tau)$) that are invisible to ordinary
homological algebra but visible in the coderived category.

The general theory of bar-cobar duality in the coderived
setting is developed in
Chapter~\ref{chap:bar-cobar}
(Theorem~\ref{thm:bar-cobar-isomorphism-main}).

% ================================================================
\section{Genus~2 --- Siegel modular forms}\label{sec:frame-genus2}
% ================================================================

At genus~$2$, the Riemann surface $\Sigma_2$ has period matrix
$\Omega = \bigl(\begin{smallmatrix} \tau_1 & z \\ z & \tau_2
\end{smallmatrix}\bigr) \in \mathcal{H}_2$ (the Siegel upper
half-space).  The genus-$2$ propagator is
\begin{equation}\label{eq:frame-genus2-propagator}
K^{(2)}(w_1, w_2 \,|\, \Omega)
= \partial_{w_1} \log E(w_1, w_2;\, \Omega)
\end{equation}
where $E(w_1, w_2;\, \Omega)$ is the prime form on~$\Sigma_2$.

Near $w_1 = w_2$, the prime form has the expansion
\begin{equation}\label{eq:frame-prime-form-expansion}
E(w_1, w_2;\, \Omega)
= (w_1 - w_2)\bigl(1 + c_2(\Omega)(w_1 - w_2)^2
+ c_4(\Omega)(w_1 - w_2)^4 + \cdots\bigr)
\end{equation}
where the coefficients $c_{2j}(\Omega)$ are Siegel modular forms
for $\mathrm{Sp}(4, \mathbb{Z})$.

Three features invisible at genus~$1$ emerge:
\begin{enumerate}
\item \emph{Multi-dimensional Hodge bundle.}
  At genus~$1$, the Hodge bundle $\mathbb{E}$ is a line bundle and
  $\lambda_1 = c_1(\mathbb{E})$.  At genus~$2$, $\mathbb{E}$ has
  rank~$2$, and $\lambda_2 = c_2(\mathbb{E})$ is a determinantal
  class involving the interaction of both B-cycle monodromies.

\item \emph{Off-diagonal periods.}
  The entry $z = \Omega_{12}$ encodes the interaction between the
  two handles of $\Sigma_2$.

\item \emph{Function-valued monodromies.}
  The genus-$1$ B-cycle monodromy of the propagator is the
  constant~$-2\pi i$.  At genus~$2$, the monodromies
  $-2\pi i \cdot \omega_j(w)$ are sections of
  $\omega_{\Sigma_2}$.  Each B-cycle defect contributes a
  $2$-form on moduli space, and the two defects combine
  multiplicatively: hence $\lambda_2$ lives in $H^4$
  (not~$H^2$).
\end{enumerate}

\begin{theorem}[Genus-2 curvature]\label{thm:frame-genus2-curvature}
The genus-$2$ bar complex of $\mathcal{H}_k$ has curvature:
\begin{equation}\label{eq:frame-genus2-curvature}
(d^{(2)})^2 = k \cdot \lambda_2
\end{equation}
where $\lambda_2 = c_2(\mathbb{E}) \in H^4(\overline{\mathcal{M}}_2)$.
\end{theorem}

The free energy at genus~$2$ is
\begin{equation}\label{eq:frame-F2}
F_2(\mathcal{H}_k) = k \cdot
\int_{\overline{\mathcal{M}}_{2,1}} \psi_1^2\, \lambda_2
= k \cdot \frac{7}{8} \cdot \frac{|B_4|}{4!}
= \frac{7k}{5760}
\end{equation}
where $B_4 = -1/30$ is the fourth Bernoulli number.

% ================================================================
\section{The genus tower --- recognizing the $\hat{A}$-genus}
\label{sec:frame-genus-tower}
% ================================================================

The pattern from genus~$1$ and genus~$2$ extends.  At genus~$g$,
the bar complex curvature is
\begin{equation}\label{eq:frame-genus-g-curvature}
(d^{(g)})^2 = k \cdot \lambda_g
\end{equation}
where $\lambda_g = c_g(\mathbb{E}) \in H^{2g}(\overline{\mathcal{M}}_g)$
is the top Chern class of the Hodge bundle
(Theorem~\ref{thm:heisenberg-obs}).

The genus-$g$ free energy is
\begin{equation}\label{eq:frame-Fg}
F_g(\mathcal{H}_k) = k \cdot
\frac{2^{2g-1}-1}{2^{2g-1}} \cdot \frac{|B_{2g}|}{(2g)!}
\end{equation}
for all $g \geq 1$ (Theorem~\ref{thm:heisenberg-all-genera}).
The first values are:
\begin{align}
F_1 &= k \cdot \frac{1}{24}, \label{eq:frame-F1} \\
F_2 &= k \cdot \frac{7}{5760}, \label{eq:frame-F2-repeat} \\
F_3 &= k \cdot \frac{31}{967680}, \label{eq:frame-F3} \\
F_4 &= k \cdot \frac{127}{154828800}. \label{eq:frame-F4}
\end{align}

The generating function assembles these into a closed form:
\begin{equation}\label{eq:frame-generating-function}
\sum_{g=1}^{\infty} F_g(\mathcal{H}_k)\, x^{2g}
= k\left(\frac{x/2}{\sin(x/2)} - 1\right).
\end{equation}

Including the leading term ($g = 0$, where $F_0 = 1$ by convention):
\begin{equation}\label{eq:frame-full-generating}
1 + \sum_{g=1}^{\infty} F_g\, x^{2g}
= \frac{x/2}{\sin(x/2)}.
\end{equation}

The function $x/2 / \sin(x/2)$ is the $\hat{A}$-genus evaluated at
imaginary argument.  The $\hat{A}$-genus of a complex manifold~$M$
is the multiplicative genus associated to the power series
\begin{equation}\label{eq:frame-a-hat-def}
\hat{A}(x) = \frac{x/2}{\sinh(x/2)}
= 1 - \frac{x^2}{24} + \frac{7x^4}{5760} - \cdots
\end{equation}
and the Wick rotation $x \mapsto ix$ sends $\sinh(x/2)$ to
$i\sin(x/2)$, producing precisely the generating
function~\eqref{eq:frame-full-generating}.

This identification is not coincidental: it follows from the
Faber--Pandharipande $\lambda_g$ formula and the Grothendieck--Riemann--Roch
theorem.  We derive both connections here.

\subsubsection*{The Faber--Pandharipande formula}

The genus-$g$ obstruction class is
$\mathrm{obs}_g(\mathcal{H}_k) = k \cdot \lambda_g$
where $\lambda_g = c_g(\mathbb{E}) \in H^{2g}(\overline{\mathcal{M}}_g)$
is the top Chern class of the Hodge bundle
$\mathbb{E} = \pi_*\omega_{\mathcal{C}/\overline{\mathcal{M}}_g}$.
The free energy extracts this as a number via integration against
powers of the cotangent line class $\psi_1 = c_1(\mathbb{L}_1)$
on $\overline{\mathcal{M}}_{g,1}$:
\begin{equation}\label{eq:frame-fp-integral}
F_g(\mathcal{H}_k) = k \cdot
\int_{\overline{\mathcal{M}}_{g,1}} \psi_1^{2g-2}\, \lambda_g.
\end{equation}
The Faber--Pandharipande formula~\cite{FP00} evaluates this integral:
\begin{equation}\label{eq:frame-fp-formula}
\int_{\overline{\mathcal{M}}_{g,1}} \psi_1^{2g-2}\, \lambda_g
= \frac{2^{2g-1}-1}{2^{2g-1}} \cdot \frac{|B_{2g}|}{(2g)!}
\end{equation}
where $B_{2g}$ is the $2g$-th Bernoulli number.

The connection to the $\hat{A}$-genus is visible at the level of
generating functions.  The Bernoulli numbers are defined by
\begin{equation}\label{eq:frame-bernoulli-gf}
\frac{x}{e^x - 1} = \sum_{n=0}^{\infty} \frac{B_n}{n!}\, x^n
= 1 - \frac{x}{2} + \frac{x^2}{12} - \frac{x^4}{720} + \cdots
\end{equation}
which is the Todd genus generating function.  The factor
$(2^{2g-1}-1)/2^{2g-1}$ in~\eqref{eq:frame-fp-formula} arises
from the relationship between the Todd class and the $\hat{A}$-class:
\begin{equation}\label{eq:frame-todd-ahat}
\mathrm{Td}(x) = \frac{x}{1 - e^{-x}} = e^{x/2} \cdot \hat{A}(x),
\qquad
\hat{A}(x) = \frac{x/2}{\sinh(x/2)}.
\end{equation}
Expanding $\hat{A}(x) = \sum_{g=0}^{\infty} a_g\, x^{2g}$
(the $\hat{A}$-class is even), the coefficients are
\begin{equation}\label{eq:frame-ahat-bernoulli}
a_g = \frac{(-1)^g(2^{2g-1}-1)}{2^{2g-1}} \cdot \frac{B_{2g}}{(2g)!}
\end{equation}
which matches~\eqref{eq:frame-fp-formula} up to the sign
$(-1)^g$ absorbed by the Wick rotation $x \mapsto ix$.
This confirms equation~\eqref{eq:frame-full-generating}.

\subsubsection*{The GRR derivation}

\begin{remark}[The family index theorem]
\label{rem:frame-family-index}
The appearance of the $\hat{A}$-genus admits a conceptual explanation
via Grothendieck--Riemann--Roch (Theorem~\ref{thm:family-index}).
For $\mathcal{A} = \mathcal{H}_k$, the modular deformation complex
is $\mathcal{D}_{\mathcal{H}_k}^{(g)} = k \cdot \mathbb{E}$
($k$ copies of the Hodge bundle).  The GRR formula applied to the
forgetful morphism $\pi \colon \overline{\mathcal{M}}_{g,1}
\to \overline{\mathcal{M}}_g$ gives
\begin{equation}\label{eq:frame-grr}
\mathrm{ch}(\pi_!\, \omega_{\mathcal{C}})
= \pi_*\!\left(\mathrm{ch}(\omega_{\mathcal{C}}) \cdot
\mathrm{Td}(T_\pi)\right)
\end{equation}
where $T_\pi$ is the relative tangent bundle.  Since
$\omega_{\mathcal{C}} = T_\pi^*$, the pushforward~\eqref{eq:frame-grr}
produces the Chern character of $\mathbb{E}$ in terms of
$\kappa$-classes, recovering the Mumford formula
$\mathrm{ch}_g(\mathbb{E}) = B_{2g} \cdot \kappa_{2g-1}/(2g)!$.
Combining with the bar-complex obstruction
$\mathrm{obs}_g = k \cdot \lambda_g$
and integrating recovers~\eqref{eq:frame-generating-function}
directly.

The genus expansion of \emph{every} Koszul chiral algebra is an index:
the $\hat{A}$-class pushforward of the Hodge bundle scaled by the
bar-complex curvature.
\end{remark}

\begin{table}[ht]
\centering
\caption{Free energy values $F_g$ for $\mathcal{H}_k$ ($k = 1$),
genera $1 \leq g \leq 6$}
\label{tab:frame-free-energy}
\begin{tabular}{|c|c|c|c|}
\hline
$g$ & $F_g$ (exact) & $F_g$ (decimal) &
  Bernoulli $|B_{2g}|$ \\
\hline
1 & $1/24$ & $4.17 \times 10^{-2}$ & $1/6$ \\
2 & $7/5760$ & $1.22 \times 10^{-3}$ & $1/30$ \\
3 & $31/967680$ & $3.20 \times 10^{-5}$ & $1/42$ \\
4 & $127/154828800$ & $8.20 \times 10^{-7}$ & $1/30$ \\
5 & $511/24524881920$ & $2.08 \times 10^{-8}$ & $5/66$ \\
6 & $1414477/2678117105664000$ & $5.28 \times 10^{-10}$ &
$691/2730$ \\
\hline
\end{tabular}
\end{table}

The generating function~\eqref{eq:frame-generating-function} has
radius of convergence $|x| < 2\pi$ (the first zero of $\sin(x/2)$
is at $x = 2\pi$).  The free energies $F_g$ decrease factorially,
and the genus expansion is convergent --- unlike the
genus expansion of the physical partition function, which grows as
$(2g)!$ (Mirzakhani volumes) and is only asymptotic.

% ================================================================
\section{Complementarity --- Theorem C in action}
\label{sec:frame-complementarity}
% ================================================================

The Heisenberg algebra $\mathcal{H}_k$ and its Koszul dual
$\mathcal{H}_k^! = \mathrm{Sym}^{\mathrm{ch}}(V^*)$ form
a Koszul pair.  Their genus-$g$ quantum corrections are
complementary.

\begin{theorem}[Quantum complementarity for Heisenberg]
\label{thm:frame-complementarity}
For the Koszul pair
$(\mathcal{H}_k, \mathrm{Sym}^{\mathrm{ch}}(V^*))$
at genus~$g$:
\begin{equation}\label{eq:frame-complementarity}
Q_g(\mathcal{H}_k) + Q_g(\mathcal{H}_k^!)
= H^*(\overline{\mathcal{M}}_g,\, \mathcal{Z}(\mathcal{H}_k))
\end{equation}
where $Q_g$ denotes the genus-$g$ quantum correction (the
obstruction to extending the genus-$0$ bar-cobar
quasi-isomorphism to genus~$g$), and $\mathcal{Z}(\mathcal{H}_k)$
is the center local system on~$\overline{\mathcal{M}}_g$.
\end{theorem}

At genus~$1$:
\begin{itemize}
\item The Heisenberg deformation space is one-dimensional,
  generated by the level~$k$.
\item The obstruction space of $\mathrm{Sym}^{\mathrm{ch}}(V^*)$
  is also one-dimensional, generated by the curvature
  $m_0 = -k \cdot \omega$.
\item Together, they span $H^*(\overline{\mathcal{M}}_1,
  \mathcal{Z}(\mathcal{H}_k)) \cong \mathbb{C}^2$.
\end{itemize}

The complementarity is controlled by Verdier duality on
$\overline{\mathcal{M}}_g$: the bar complex computes
deformations, the cobar complex computes obstructions,
and Verdier duality exchanges them.  What $\mathcal{H}_k$ sees
as a deformation (varying the level~$k$), its Koszul dual
sees as an obstruction (the curvature $m_0$), and vice versa.

This is Theorem~C
(Theorem~\ref{thm:quantum-complementarity-main})
for the Heisenberg case.  The general theorem establishes
complementarity for all Koszul chiral algebra pairs, with the
ambient space $H^*(\overline{\mathcal{M}}_g, \mathcal{Z})$
carrying a shifted-symplectic pairing from Verdier duality on
the center local system.

% ================================================================
\section{The modular characteristic package}
\label{sec:frame-modular-package}
% ================================================================

Everything computed in this chapter --- the bar complex,
the Koszul dual, the curvature at each genus, the free energy
tower, the complementarity --- assembles into a single object.

\begin{definition}[Modular characteristic package for Heisenberg]
\label{def:frame-modular-package}
The modular characteristic package of $\mathcal{H}_k$ consists of:
\begin{enumerate}
\item The \emph{universal Maurer--Cartan class}
  \begin{equation}\label{eq:frame-MC-class}
  \Theta_{\mathcal{H}_k} \in
  \mathrm{MC}\!\left(
    \mathrm{Def}_{\mathrm{cyc}}(\mathcal{H}_k)
    \;\widehat{\otimes}\;
    R\Gamma\!\left(\overline{\mathcal{M}}_{g,\bullet},\,
    \mathbb{Q}\right)
  \right)
  \end{equation}
  whose image in the scalar center produces the obstruction
  coefficients $\mathrm{obs}_g = k \cdot \lambda_g$.

\item The \emph{ambient deformation complex}
  \begin{equation}\label{eq:frame-ambient}
  H_{\mathcal{H}_k}
  := R\Gamma(\overline{\mathcal{M}}_g,\,
  \mathcal{Z}_{\mathcal{H}_k})
  \end{equation}
  with Verdier duality providing the shifted-symplectic structure
  for complementarity.

\item The \emph{spectral discriminant}
  $\Delta_{\mathcal{H}_k}(x) = 1$
  (trivial for rank~$1$; for affine Kac--Moody algebras, the
  discriminant acquires the form
  $\Delta(x) = (1-3x)(1+x)$).

\item The \emph{obstruction coefficient}
  $\kappa(\mathcal{H}_k) = k$, the first characteristic number
  of~$\Theta_{\mathcal{H}_k}$.
\end{enumerate}
\end{definition}

For the Heisenberg algebra, the modular characteristic package is
fully determined by the single number $\kappa = k$.  This is an
artifact of the abelian (rank-$1$) structure.  In
Chapters~\ref{chap:kac-moody-koszul}
and~\ref{chap:w-algebras}, the
package acquires non-trivial spectral discriminant, non-scalar
Maurer--Cartan class, and richer ambient deformation complex.

\begin{remark}[The first characteristic number is not the
  fundamental object]\label{rem:frame-kappa-not-fundamental}
The genus-$1$ obstruction coefficient $\kappa(\mathcal{A})$ ---
which equals the level~$k$ for Heisenberg, the Sugawara central
charge $c/2$ for Virasoro, and $(k+h^\vee)\dim(\mathfrak{g})/(2h^\vee)$
for affine Kac--Moody --- is often treated as the primary invariant.
It is not.  The fundamental object is the universal
class~$\Theta_{\mathcal{A}}$, of which $\kappa$ is merely the
first characteristic number.  The complete genus tower, the
complementarity pairing, and the index theorem all follow from
$\Theta_{\mathcal{A}}$, not from $\kappa$ alone.
\end{remark}

% ================================================================
\section{What breaks for $\widehat{\mathfrak{sl}}_2$ ---
  preview of the non-abelian theory}
\label{sec:frame-nonabelian-preview}
% ================================================================

The Heisenberg algebra has a single current with a double-pole
self-OPE and no simple pole.  The simplest non-abelian chiral
algebra, $\widehat{\mathfrak{sl}}_2$ at level~$k$, has three
currents $e(z), f(z), h(z)$ with OPE:
\begin{align}
e(z)\, f(w) &\sim \frac{h(w)}{z-w}
  + \frac{k}{(z-w)^2}, \label{eq:frame-sl2-ef} \\
h(z)\, e(w) &\sim \frac{2e(w)}{z-w}, \qquad
h(z)\, f(w) \sim \frac{-2f(w)}{z-w},
  \label{eq:frame-sl2-he-hf} \\
h(z)\, h(w) &\sim \frac{2k}{(z-w)^2}.
  \label{eq:frame-sl2-hh}
\end{align}

The simple poles in~\eqref{eq:frame-sl2-ef}--\eqref{eq:frame-sl2-he-hf}
are the new feature.  They contribute a
\emph{bracket component} $d_{\mathrm{bracket}}$ to the bar
differential, which was identically zero for Heisenberg.

\subsection{The bar differential at degrees $0$--$2$}
\label{subsec:frame-sl2-bar-low-degree}

We repeat the Heisenberg computation of \S\ref{sec:frame-bar-deg1}--\S\ref{sec:frame-bar-deg2} for $\widehat{\mathfrak{sl}}_2$, emphasizing what changes.

\subsubsection*{Degree $0$}
The degree-$0$ bar space is $\bar{B}_0^{\mathrm{ch}} = \mathbb{C}$,
as before: the augmentation ideal modulo products.

\subsubsection*{Degree $1$: two types of pole}
A degree-$1$ element now involves two currents from $\{e, f, h\}$
and a single logarithmic form $\eta_{12}$.  The bar differential
$d = d_{\mathrm{bracket}} + d_{\mathrm{residue}}$ has two
components, one for each pole order.

\medskip\noindent\textbf{Simple-pole component $d_{\mathrm{bracket}}$.}
The OPE $h(z)\,e(w) \sim 2e(w)/(z-w)$ has a simple pole with
residue $c_1(w) = 2e(w)$.  The bar differential extracts this as
\begin{equation}\label{eq:frame-sl2-dbracket-deg1}
d_{\mathrm{bracket}}\bigl(
  s^{-1}h(z_1) \otimes s^{-1}e(z_2) \otimes \eta_{12}
\bigr)
= c_1(z_2) = 2e(z_2).
\end{equation}
This is the Lie bracket $[h, e] = 2e$ encoded in OPE language.
Contrast with Heisenberg, where $c_1 = 0$ and $d_{\mathrm{bracket}}$
vanishes identically.  For the full $\mathfrak{sl}_2$ system,
$d_{\mathrm{bracket}}$ on degree-$1$ elements recovers the
Chevalley--Eilenberg differential of $\mathfrak{sl}_2$:
\begin{equation}\label{eq:frame-sl2-ce-table}
\begin{aligned}
d_{\mathrm{bracket}}(h \otimes e \otimes \eta) &= 2e, &\qquad
d_{\mathrm{bracket}}(h \otimes f \otimes \eta) &= -2f, \\
d_{\mathrm{bracket}}(e \otimes f \otimes \eta) &= h, &\qquad
d_{\mathrm{bracket}}(h \otimes h \otimes \eta) &= 0.
\end{aligned}
\end{equation}

\medskip\noindent\textbf{Double-pole component $d_{\mathrm{residue}}$.}
The double poles $h(z)\,h(w) \sim 2k/(z-w)^2$ and
$e(z)\,f(w) \sim k/(z-w)^2 + \cdots$ contribute a curvature component:
\begin{equation}\label{eq:frame-sl2-dresid-deg1}
\begin{aligned}
d_{\mathrm{residue}}(h \otimes h \otimes \eta) &= 2k, \\
d_{\mathrm{residue}}(e \otimes f \otimes \eta) &= k, \\
d_{\mathrm{residue}}(h \otimes e \otimes \eta) &= 0.
\end{aligned}
\end{equation}
This is the Killing form $\kappa(J^a, J^b) = k \cdot (J^a, J^b)$
in OPE language, identical in structure to the Heisenberg double-pole
extraction.

The full bar differential is $d = d_{\mathrm{bracket}} + d_{\mathrm{residue}}$.
At degree~$1$:
\begin{equation}\label{eq:frame-sl2-full-d-deg1}
d(e \otimes f \otimes \eta) = \underbrace{h}_{\text{bracket}} + \underbrace{k}_{\text{curvature}}, \qquad
d(h \otimes h \otimes \eta) = \underbrace{0}_{\text{bracket}} + \underbrace{2k}_{\text{curvature}}.
\end{equation}

\subsubsection*{Degree $2$: the Jacobi failure}

At degree~$2$, we have three currents and two logarithmic forms
$\eta_{12} \wedge \eta_{23}$.  The bracket component
$d_{\mathrm{bracket}}$ now applies a Lie bracket to two of the
three currents via a simple-pole extraction, producing a degree-$1$
element.  Consider:
\begin{equation}\label{eq:frame-sl2-deg2-element}
\xi = e(z_1) \otimes h(z_2) \otimes f(z_3)
  \otimes \eta_{12} \wedge \eta_{23}.
\end{equation}
The bracket component extracts simple poles at each collision:
\begin{align}
\partial_{D_{12}}^{\mathrm{bracket}}(\xi)
  &= [e, h](z_{12}) \otimes f(z_3) \otimes \eta_{23}
  = -2e(z_{12}) \otimes f(z_3) \otimes \eta_{23},
\label{eq:frame-sl2-D12-bracket} \\
\partial_{D_{23}}^{\mathrm{bracket}}(\xi)
  &= -e(z_1) \otimes [h, f](z_{23}) \otimes \eta_{12}
  = +2e(z_1) \otimes f(z_{23}) \otimes \eta_{12},
\label{eq:frame-sl2-D23-bracket}
\end{align}
where the sign in~\eqref{eq:frame-sl2-D23-bracket} comes from the
wedge reordering $\eta_{12} \wedge \eta_{23}
= -\eta_{23} \wedge \eta_{12}$, as in the Heisenberg computation
of \S\ref{sec:frame-bar-deg2}.

Now apply $d_{\mathrm{bracket}}$ again.  From~\eqref{eq:frame-sl2-D12-bracket}:
\[
d_{\mathrm{bracket}}\bigl(-2e \otimes f \otimes \eta_{23}\bigr)
= -2h.
\]
From~\eqref{eq:frame-sl2-D23-bracket}:
\[
d_{\mathrm{bracket}}\bigl(2e \otimes f \otimes \eta_{12}\bigr)
= 2h.
\]
These would cancel if $d_{\mathrm{bracket}}$ were a differential.
But we have neglected the $D_{13}$ stratum (the $e$--$f$ collision),
which contributes via $[e, f] = h$:
\begin{equation}\label{eq:frame-sl2-D13-bracket}
\partial_{D_{13}}^{\mathrm{bracket}}(\xi)
= h(z_{13}) \otimes h(z_2) \otimes \eta_{12}|_{D_{13}}.
\end{equation}
After careful sign accounting (the form $\eta_{12}\wedge\eta_{23}$
restricts to $D_{13}$ non-trivially, unlike the Heisenberg case),
the three terms of $d_{\mathrm{bracket}}^2$ on $\xi$ do \emph{not}
cancel.  The residual term is proportional to $k$: it is precisely
the double-pole contribution from the $e$--$f$ OPE that the
bracket component ignores.

\begin{remark}[Why $d_{\mathrm{bracket}}^2 \neq 0$]
\label{rem:frame-why-dbracket-fails}
In the classical bar complex of $U(\mathfrak{g})$, the
Chevalley--Eilenberg differential satisfies
$d_{\mathrm{CE}}^2 = 0$ by the Jacobi identity.  In the chiral
setting, $d_{\mathrm{bracket}}$ extracts simple poles, but the
Borcherds identity couples simple and double poles:
the ``Jacobi'' relation $[e, [h, f]] - [h, [e, f]] - [[e,h], f] = 0$
holds only when the double-pole (level-$k$) terms are included.
The full differential $d = d_{\mathrm{bracket}} + d_{\mathrm{residue}}$
satisfies $d^2 = 0$ via the complete Borcherds identity, which
intertwines both pole orders.  This is the chiral Koszul phenomenon:
the two components of $d$ are individually \emph{not} differentials;
they form a differential only in combination.
\end{remark}

\subsection{$d_{\mathrm{bracket}}^2 \neq 0$}

For $\widehat{\mathfrak{sl}}_2$, the bracket component
$d_{\mathrm{bracket}}$ alone does \emph{not} satisfy
$d_{\mathrm{bracket}}^2 = 0$.  This is proved computationally:
all $2048$ sign conventions for the bar differential have been
verified to give $d_{\mathrm{bracket}}^2 \neq 0$
(see Theorem~\ref{thm:pbw-koszulness-criterion}).
The combined differential $d = d_{\mathrm{bracket}}
+ d_{\mathrm{curvature}}$ satisfies $d^2 = 0$ by the full
Borcherds identity (which involves both simple and double poles),
not by the Jacobi identity alone.

This is the essential difference between classical and chiral
Koszul duality.  In the classical setting, the enveloping algebra
$U(\mathfrak{g})$ is Koszul (BGS~\cite{BGS96}), and the bar
differential $d_{\mathrm{CE}} = d_{\mathrm{bracket}}$ satisfies
$d_{\mathrm{CE}}^2 = 0$ by the Jacobi identity.  In the chiral
setting, $\widehat{\mathfrak{g}}_k$ uses all OPE poles (the full
Borcherds identity), and the bracket-only piece is not a
differential on its own.

\subsection{The PBW spectral sequence}

The correct computational tool is the \emph{PBW spectral sequence}.
Filter the bar complex by conformal weight.  The associated graded
is the classical bar complex of $U(\mathfrak{g})$, which \emph{is}
a chain complex (with differential $d_{\mathrm{CE}}$).  The
spectral sequence has:
\begin{itemize}
\item $E_1$ page: the bar complex of $U(\mathfrak{g})$ with
  $d_{\mathrm{CE}}$.
\item $E_2$ page: the Chevalley--Eilenberg cohomology
  $H^*_{\mathrm{CE}}(\mathfrak{g}, \mathbb{C})$.
\end{itemize}
For Koszul algebras, the spectral sequence collapses at~$E_2$,
giving the chiral bar cohomology.  The proof that
$\widehat{\mathfrak{g}}_k$ is chiral Koszul uses PBW flatness
of the associated graded together with the classical Koszulness
of $U(\mathfrak{g})$ (Priddy~\cite{Priddy70}); see
Theorem~\ref{thm:pbw-koszulness-criterion}.

\subsection{The Feigin--Frenkel involution as Koszul duality}

The Koszul dual of $\widehat{\mathfrak{g}}_k$ is
$\widehat{\mathfrak{g}}_{-k-2h^\vee}$, where $h^\vee$ is the
dual Coxeter number (Theorem~\ref{thm:universal-kac-moody-koszul}).
For $\mathfrak{g} = \mathfrak{sl}_2$, $h^\vee = 2$ and the duality
sends $k \mapsto -k - 4$.

\begin{center}
\small
\begin{tabular}{c|c|c}
& \textbf{Heisenberg} $\mathcal{H}_k$ &
\textbf{Kac--Moody} $\widehat{\mathfrak{g}}_k$ \\
\hline
Bracket & Abelian ($d_{\mathrm{bracket}} = 0$)
  & Non-abelian ($d_{\mathrm{bracket}} \neq 0$) \\
$d^2 = 0$ from & Arnold alone
  & Full Borcherds identity \\
Spectral sequence & Trivial ($E_1 = E_\infty$)
  & PBW ($E_2$ collapse for Koszul) \\
Koszul dual & $\mathrm{Sym}^{\mathrm{ch}}(V^*)$
  & $\widehat{\mathfrak{g}}_{-k-2h^\vee}$ \\
Level shift & None ($h^\vee = 0$)
  & $k \mapsto -k - 2h^\vee$ \\
Discriminant & $\Delta(x) = 1$
  & $\Delta(x) = (1-3x)(1+x)$ for $\mathfrak{sl}_2$ \\
Critical level & $k = 0$ (degenerate)
  & $k = -h^\vee$ (curvature vanishes)
\end{tabular}
\end{center}

At the critical level $k = -h^\vee$, the curvature of the bar
complex vanishes: $d^2 = 0$ at all genera.  This is the level at
which the center of $\widehat{\mathfrak{g}}_k$ becomes large
(the Feigin--Frenkel center $\mathfrak{z}(\widehat{\mathfrak{g}}_{-h^\vee})
\cong \mathcal{W}(\widehat{\mathfrak{g}}^\vee)$), and the bar
complex recovers the Langlands dual.

The Feigin--Frenkel involution $k \leftrightarrow -k - 2h^\vee$
is not a representation-theoretic accident: it is Koszul duality
in the chiral algebra sense, proved in
Theorem~\ref{thm:universal-kac-moody-koszul} via Serre duality
on Fulton--MacPherson configuration spaces.

% ================================================================
\section{What breaks for Yangians --- preview of the $E_1$ world}
\label{sec:frame-yangian-preview}
% ================================================================

The Heisenberg algebra is $E_\infty$-chiral: its operations are
fully commutative and local.  Configuration spaces
$C_n(X)$ are unordered.  Verdier duality on $\overline{C}_n(X)$
exchanges bar and cobar.

The Yangian $Y(\mathfrak{g})$ is $E_1$-chiral: its operations are
associative but \emph{not} commutative.  Configuration spaces are
\emph{ordered}: $\overline{C}_n^{\mathrm{ord}}(X)$ replaces
$\overline{C}_n(X)$.  The bar complex uses ordered tuples, and
the coproduct is \emph{not} cocommutative.

\begin{center}
\small
\begin{tabular}{c|c|c}
& \textbf{Heisenberg}
  ($E_\infty$-chiral) &
\textbf{Yangian} $Y(\mathfrak{g})$
  ($E_1$-chiral) \\
\hline
Configuration space & $C_n(X)$ (unordered)
  & $C_n^{\mathrm{ord}}(X)$ (ordered) \\
Coproduct & Cocommutative
  & Coassociative (not cocommutative) \\
Koszul dual & $\mathrm{Sym}^{\mathrm{ch}}(V^*)$
  & $Y_{R^{-1}}(\mathfrak{g})$ \\
Duality type & Lie $\leftrightarrow$ Com
  & $R \leftrightarrow R^{-1}$ \\
Self-dual? & No & Yes (classically)
\end{tabular}
\end{center}

The Yangian $Y(\mathfrak{g})$ is presented by the RTT relations,
which are quadratic in the $E_1$-chiral sense.  The quadratic
dual relations are the RTT relations with the inverse $R$-matrix
$R^{-1}(u)$ replacing~$R(u)$
(Theorem~\ref{thm:yangian-koszul-dual}).  At the classical limit
$\hbar = 0$, $R^{-1} = R$ and the Yangian is self-dual
(Corollary~\ref{cor:yangian-classical-self-dual}); at $\hbar \neq 0$,
the duality is a non-trivial involution.

The derived Drinfeld--Kohno theorem
(Theorems~\ref{thm:derived-dk-affine}
and~\ref{thm:derived-dk-yangian}) proves at the chain level:
\begin{equation}\label{eq:frame-derived-dk}
D^b(\mathrm{Rep}_{\mathrm{fd}}(Y_\hbar(\mathfrak{g})))
\;\simeq\;
D^b(\mathrm{Rep}_{\mathrm{fd}}(Y_{-\hbar}(\mathfrak{g})))
\end{equation}
with braiding reversal $R \mapsto R^{-1}$.  Under the Kazhdan
equivalence, this becomes $q \mapsto q^{-1}$ on quantum group
representations.  The factorization-level statement is proved on
the evaluation locus (Theorem~\ref{thm:factorization-dk-eval});
the full factorization categorical equivalence
$\mathrm{Fact}_{E_1}(Y(\mathfrak{g}))
\simeq \mathrm{Fact}_{E_1}(U_q(\mathfrak{g}))^{\mathrm{op}}$
remains conjectural
(Conjecture~\ref{conj:derived-drinfeld-kohno}).

The passage from $E_\infty$ (Heisenberg) to $E_1$ (Yangian) is
the passage from curves to higher-dimensional manifolds: $E_n$-chiral
algebras are the factorization algebras on
$\mathbb{R}^n$-stratified spaces, and the bar-cobar adjunction
generalizes to $E_n$-Koszul duality
(Theorem~\ref{thm:e1-chiral-koszul-duality}).  Chapter~\ref{ch:en-koszul-duality}
develops this framework.

% ================================================================
\section{The view from here --- what the theory wants to become}
\label{sec:frame-programme}
% ================================================================

In a single example, we have seen:
\begin{enumerate}
\item The bar complex built from configuration space integrals,
  with $d^2 = 0$ forced by the Arnold relation
  (\S\ref{sec:frame-bar-deg2}).
\item The Koszul dual read off from the bar complex cohomology:
  $\mathcal{H}_k^! = \mathrm{Sym}^{\mathrm{ch}}(V^*)$
  (\S\ref{sec:frame-koszul-dual}).
\item Bar-cobar inversion: $\Omega(\bar{B}(\mathcal{H}_k))
  \simeq \mathcal{H}_k$ (\S\ref{sec:frame-inversion}).
\item Curvature at genus~$1$: $d^2 = k \cdot \omega_1$,
  forcing the passage to the coderived category
  (\S\ref{sec:frame-genus1}).
\item The genus tower of obstruction classes
  $\mathrm{obs}_g = k \cdot \lambda_g$, with generating function
  the $\hat{A}$-genus (\S\ref{sec:frame-genus-tower}).
\item Complementarity of quantum corrections for the Koszul pair
  (\S\ref{sec:frame-complementarity}).
\end{enumerate}

These are the main theorems of this monograph, instantiated
for the Heisenberg algebra:

\begin{itemize}
\item \textbf{Theorem~A}
  (Theorem~\ref{thm:bar-cobar-isomorphism-main}):
  \emph{Geometric bar-cobar duality.}
  What we saw as ``the bar complex produces a commutative Koszul
  dual'' generalizes to: the bar functor on augmented chiral
  algebras is intertwined with Verdier duality on
  $\mathrm{Ran}(X)$, and produces the Koszul dual coalgebra
  functorially over $\overline{\mathcal{M}}_{g,n}$.

\item \textbf{Theorem~B}
  (Theorem~\ref{thm:higher-genus-inversion}):
  \emph{Bar-cobar inversion.}
  What we saw as ``cobar recovers Heisenberg'' generalizes to:
  on the Koszul locus (where the PBW spectral sequence collapses),
  the counit $\Omega \bar{B}(\mathcal{A}) \to \mathcal{A}$ is a
  quasi-isomorphism.  Off this locus, the counit persists as a
  curved quasi-isomorphism in the coderived category.

\item \textbf{Theorem~C}
  (Theorem~\ref{thm:quantum-complementarity-main}):
  \emph{Deformation-obstruction complementarity.}
  What we saw as ``deformation of $\mathcal{H}_k$ and obstruction
  of $\mathrm{Sym}^{\mathrm{ch}}(V^*)$ are complementary''
  generalizes to: for any Koszul pair $(\mathcal{A},
  \mathcal{A}^!)$, the deformation and obstruction spaces embed
  as complementary Lagrangians in a shifted-symplectic ambient
  space on $\overline{\mathcal{M}}_g$.
\end{itemize}

\subsection{The modular Koszul programme}

The four theorems, the genus tower, and the modular characteristic
package point toward a unified theory.  The \emph{modular Koszul
programme} articulates five target theorems:

\begin{enumerate}
\item[$A_{\mathrm{mod}}$.]
  \emph{Verdier-intertwined bar-cobar, functorial over
  $\overline{\mathcal{M}}_{g,n}$.}
  The bar-cobar adjunction extends to a modular operad map,
  with Verdier duality exchanging bar and cobar at every genus.

\item[$B_{\mathrm{mod}}$.]
  \emph{Inversion on the Koszul locus, coderived persistence off it.}
  The bar-cobar quasi-isomorphism holds on the Koszul locus at
  every genus; off this locus, the bar complex is curved and lives
  in the coderived category.

\item[$C_{\mathrm{mod}}$.]
  \emph{$(-1)$-shifted symplectic complementarity.}
  The deformation and obstruction spaces for a Koszul pair form
  complementary Lagrangians in a $(-1)$-shifted symplectic
  structure on $R\Gamma(\overline{\mathcal{M}}_g,
  \mathcal{Z}_{\mathcal{A}})$.

\item[\textbf{Index}.]
  \emph{Grothendieck--Riemann--Roch for modular deformations.}
  The genus series $\sum_g F_g x^{2g}$ is the $\hat{A}$-class
  pushforward of the universal Maurer--Cartan class along the
  family $\overline{\mathcal{M}}_{g,\bullet} \to \mathrm{pt}$.

\item[\textbf{DK}.]
  \emph{Derived Drinfeld--Kohno.}
  The $E_1$-factorization homology of the Yangian $Y(\mathfrak{g})$
  recovers the quantum group $U_q(\mathfrak{g})$ via a derived
  equivalence of braided tensor categories.
\end{enumerate}

Theorems~A--C are proved in this monograph
(Chapters~\ref{chap:bar-cobar}--\ref{chap:higher-genus}).
The Index theorem is proved (Theorem~\ref{thm:family-index});
Derived Drinfeld--Kohno is proved at the chain level
(Theorem~\ref{thm:derived-dk-yangian}), with the full
factorization statement conjectured.
The full programme is stated in
\S\ref{sec:modular-koszul-programme}.

\subsection{The operadic architecture}
\label{subsec:frame-operadic-architecture}

The genus tower of \S\ref{sec:frame-genus-tower} is organized by an
operadic structure that we now make visible in the Heisenberg example.

\subsubsection*{Genus $0$: $\mathrm{HyCom} \leftrightarrow \mathrm{Grav}$}

The Arnold relation (Proposition~\ref{prop:frame-arnold}) is the
defining relation of the \emph{gravity operad} $\mathrm{Grav}$,
which governs the genus-$0$ bar complex.  Its Koszul dual is the
\emph{hypercommutative operad} $\mathrm{HyCom}$, which governs the
cobar complex.  The genus-$0$ bar-cobar adjunction
\[
\bar{B} \colon \mathsf{Alg}_{\mathrm{HyCom}}
\rightleftarrows
\mathsf{coAlg}_{\mathrm{Grav}} :\! \Omega
\]
is an instance of operadic Koszul duality (Getzler--Kapranov~\cite{GK94}).
For Heisenberg:
\begin{itemize}
\item $\mathcal{H}_k$ is a $\mathrm{HyCom}$-algebra (the chiral
  operations are commutative and local).
\item $\bar{B}(\mathcal{H}_k)$ is a $\mathrm{Grav}$-coalgebra
  (the bar differential is controlled by the Arnold relation).
\item The cobar $\Omega\bar{B}(\mathcal{H}_k) \simeq \mathcal{H}_k$
  recovers the original algebra
  (\S\ref{sec:frame-inversion}).
\end{itemize}

The two operads are dual incarnations of a single structure:
$\mathrm{HyCom}(n) = H^*(\overline{M}_{0,n+1})$ (cohomology of
genus-$0$ moduli), and $\mathrm{Grav}(n) = H^*(\overline{M}_{0,n+1})^!$
(Koszul dual grading).  The genus-$0$ bar complex \emph{is} the
$\mathrm{Grav}$-coalgebra structure, and the bar differential $d^2 = 0$
follows from the $\mathrm{Grav}$ relations (= Arnold).

\subsubsection*{Genus $1$: the elliptic operad}

At genus~$1$, the moduli space $\overline{\mathcal{M}}_{1,n}$
enters.  The bar complex on an elliptic curve acquires curvature
$d^2 = k \cdot \omega_1$ (\S\ref{sec:frame-genus1}).  Operadically,
the genus-$1$ structure is governed by the \emph{elliptic operad}:
the collection $\{H^*(\overline{\mathcal{M}}_{1,n})\}_{n \geq 1}$
with composition maps from clutching morphisms.

For Heisenberg, the elliptic operad acts as follows.  The
$E_2$ non-holomorphicity (\S\ref{sec:frame-holomorphic-anomaly})
is the genus-$1$ modular completion of the genus-$0$ Arnold
relation: the Eisenstein series $E_2(\tau)$ fails to be modular
precisely because the ``boundary of boundary'' at genus~$1$
acquires the nodal contribution $\delta_0 \in
\partial\overline{\mathcal{M}}_{1,1}$.  This is the clutching
morphism $\overline{\mathcal{M}}_{0,3} \to
\partial\overline{\mathcal{M}}_{1,1}$ that glues a genus-$0$
curve to itself.

\subsubsection*{All genera: the modular operad and clutching}

The full genus tower is organized by the \emph{modular operad}
$\{\overline{\mathcal{M}}_{g,n}\}_{2g-2+n>0}$
(Getzler--Kapranov~\cite{GK98}).  A modular operad has two types of
composition:
\begin{enumerate}
\item \emph{Tree-level composition} ($\circ_i$): gluing marked
  points on different curves. This is the genus-$0$ operadic
  structure.
\item \emph{Self-gluing} ($\xi_{ij}$): identifying two marked
  points on the \emph{same} curve, increasing genus by~$1$.
  The clutching morphism
  $\xi \colon \overline{\mathcal{M}}_{g,n+2}
  \to \overline{\mathcal{M}}_{g+1,n}$
  glues marked points $i$ and $j$.
\end{enumerate}

For Heisenberg, self-gluing is the mechanism that produces the
genus tower:
\begin{itemize}
\item At genus~$0$: tree-level composition gives $d^2 = 0$
  (Arnold).
\item At genus~$1$: self-gluing of a genus-$0$ curve introduces
  $d^2 = k \cdot \omega_1$ (curvature from the new cycle).
\item At genus~$g$: iterated self-gluing produces
  $\mathrm{obs}_g = k \cdot \lambda_g$, with $\lambda_g$ emerging
  from $g$ independent cycles on the nodal degeneration.
\end{itemize}

The \emph{Feynman transform} of a modular operad (Getzler--Kapranov~\cite{GK98})
provides the all-genus bar construction.  For the modular operad
$\overline{\mathcal{M}}$, the Feynman transform produces a twisted
modular operad whose algebras are precisely curved
$A_\infty$-algebras with the curvature controlled by the
self-gluing maps.  This is why the universal Maurer--Cartan class
$\Theta_{\mathcal{A}}$ of Definition~\ref{def:frame-modular-package}
lives in $\mathrm{Def}_{\mathrm{cyc}}(\mathcal{A})
\,\widehat{\otimes}\, R\Gamma(\overline{\mathcal{M}}_{g,\bullet},
\mathbb{Q})$: the cyclic deformations of $\mathcal{A}$ are coupled
to the cohomology of moduli by the Feynman transform, and the
genus-$g$ component of $\Theta_{\mathcal{A}}$ encodes how the
self-gluing maps act on the bar-complex curvature.

\subsection{The Heisenberg algebra as atom}

The Heisenberg algebra is the atom from which this theory is built.
Every theorem in this monograph is an elaboration of what we have
seen in this chapter:
\begin{itemize}
\item The Arnold relation is the genus-$0$ seed of all bar complex
  constructions.
\item Verdier duality on $\mathrm{Ran}(X)$ makes the Koszul dual
  a theorem, not a definition.
\item The genus-$1$ curvature $d^2 = \kappa \cdot \omega_1$ forces
  the coderived category.
\item The clutching of stable curves (modular operad structure)
  organizes the genus tower.
\end{itemize}

These are the four irreducible pieces of the theory: Arnold
relation, Verdier duality, genus-$1$ curvature, and clutching.
The remainder of this monograph develops them in full generality,
applies them to every major family of chiral algebras, and
connects them to the physics of conformal field theory,
string theory, and gauge theory.
